% !TEX program = tectonic
\documentclass[11pt,reqno]{amsart}
\usepackage[margin=1.15in]{geometry}
\usepackage{amsmath,amssymb,amsthm,mathtools}
\usepackage{booktabs}
\usepackage[colorlinks=true,linkcolor=blue,citecolor=blue]{hyperref}

\newtheorem{theorem}{Theorem}[section]
\newtheorem{proposition}[theorem]{Proposition}
\newtheorem{lemma}[theorem]{Lemma}
\newtheorem{corollary}[theorem]{Corollary}
\theoremstyle{definition}
\newtheorem{conjecture}[theorem]{Conjecture}
\newtheorem{remark}[theorem]{Remark}
\newtheorem{definition}[theorem]{Definition}

\DeclareMathOperator{\Tr}{Tr}
\DeclareMathOperator{\tr}{tr}
\DeclareMathOperator{\HS}{HS}
\DeclareMathOperator{\Nm}{N}
\newcommand{\FF}{\mathbb{F}}
\newcommand{\ZZ}{\mathbb{Z}}
\newcommand{\CC}{\mathbb{C}}
\newcommand{\RR}{\mathbb{R}}
\newcommand{\QQ}{\mathbb{Q}}
\newcommand{\ep}{\mathrm{e}_p}
\newcommand{\HH}{\mathcal{H}}
\newcommand{\GUE}{\mathrm{GUE}}
\newcommand{\GOE}{\mathrm{GOE}}
\newcommand{\Kl}{\mathrm{Kl}}

\begin{document}

\title[Exact Wigner--Dyson statistics and subgroup Kloosterman moments]
{Exact Wigner--Dyson statistics for Gaussian Weyl observables on quantised cat
maps, and exact moments of subgroup Kloosterman sums}

\author{Churilov M V\\\texttt{churilovm1305@gmail.com}}
\date{30 August 2026}

\begin{abstract}
We prove two exact statements bearing on the Bohigas--Giannoni--Schmit (BGS)
problem for quantised linear automorphisms of the torus.

First, let $U_A$ be the Weyl quantisation of a hyperbolic $A\in SL_2(\ZZ)$ at
odd prime level $N$, let $\HH_\lambda$ be an eigenspace of $U_A$ of dimension
$m$, and let $V=\sum_{u\neq 0}g_u\,T(u)$ be the Gaussian Weyl observable
($g_{-u}=\overline{g_u}$, otherwise i.i.d.\ standard complex Gaussians). We show
that the compression $P_\lambda VP_\lambda$ has, \emph{exactly} and not merely
asymptotically, the law of $\sqrt{N}\,\GUE_m$ modulo the one-dimensional trace
direction; consequently every statistic of the eigenvalue differences of the
first-order multiplet splitting coincides exactly with the corresponding
$\GUE_m$ statistic. If $V$ is projected onto the subspace invariant under the
arithmetic time-reversal $\Theta$, the same compression becomes exactly
$\sqrt{N/2}\,\GOE_m$ in the $\Theta$-real basis. The only input is the
tight-frame identity $\sum_u T(u)XT(u)^*=N\Tr(X)\,\mathbf 1$; in particular the
statement holds for an \emph{arbitrary} $m$-dimensional subspace, so it carries
no arithmetic information, and no numerical experiment performed with a random
Weyl observable can be evidence for the arithmetic content of BGS.

Second, we determine exactly all moments of the subgroup Kloosterman sums
$S_H(a,b)=\sum_{h\in H}\ep(ah+bh^{-1})$, $H\le\FF_p^\times$ of prime order $k$,
in the range $p>(u+v)^{k-1}$: the mixed moments vanish identically and the
balanced moments equal $p^2\,\mathrm{P}(s,k)$ with
$\mathrm{P}(s,k)=s!^2[x^s]I_0(2\sqrt x)^k$, the exact complex-Gaussian value up
to an explicit finite-$k$ correction. The proof replaces the mod-$p$ count by a
vanishing-sum count in $\ZZ[\zeta_k]$ and uses that a balanced $\pm1$ vanishing
sum of $k$-th roots of unity is trivial when $k$ is prime. We deduce a
sub-Gaussian tail bound and a central limit theorem, and we contrast this with
the Sato--Tate law at $k=p-1$. We further prove three exact structural theorems. Theorem C: the
antiunitary square is $+1$ unconditionally, because $M=SR$ always has
equal diagonal entries, which closes the standing problem O6 and removes
a hypothesis. Theorem D: the parity operator commutes with $U_A$, and it
acts as a scalar on eigenspaces exactly when $A^k=-\mathbf1$; otherwise
every observable with a real (even) Weyl symbol has a block-diagonal
compression and therefore cannot be Wigner--Dyson. Theorem E: the Hecke
torus supplies, off a locus of density $O(1/p)$, an antiunitary that does
commute with the single-orbit compression, forcing the orthogonal class.
After desymmetrising by parity, the fully deterministic single-observable
compression is indistinguishable from $\GOE$ at $m$ up to $337$ per block
by an unfolding-free ratio statistic, while a matched arithmetic-free
surrogate is indistinguishable from $\GUE$: BGS holds numerically, with
the symmetry class correctly predicted. All statements are certified
numerically at machine precision; the exceptional prime sets are computed
exhaustively. Finally we construct the Hecke eigenbasis correctly and find a sharp dichotomy in it: the diagonal matrix elements of a pure Weyl operator have excess kurtosis $-1.00\pm0.02$, exactly the Sato--Tate value and $22$--$65\sigma$ from Gaussian, while the off-diagonal elements are Gaussian to four decimals ($-0.0014\pm0.0019$ over $5.8\cdot10^6$ samples). The arithmetic non-Gaussianity is confined to the diagonal, where it cannot create or destroy level repulsion.
\end{abstract}

\maketitle

\section{Introduction}

The Bohigas--Giannoni--Schmit conjecture \cite{BGS1984} asserts that the
spectral fluctuation statistics of a quantised classically chaotic system agree,
in the semiclassical limit, with those of the random-matrix ensemble selected by
the system's antiunitary symmetries. No instance of the conjecture is known for
a single deterministic system; the strongest semiclassical results establish the
form factor to all orders in perturbation theory
\cite{SieberRichter2001,Mueller2004} rather than as a convergent statement.

Quantised linear automorphisms of the torus (``cat maps'') occupy a special
place. They are exactly solvable, and precisely for that reason they are known
\emph{not} to satisfy BGS in its naive form: the propagator satisfies
$U_A^{K}=\chi\cdot\mathbf 1$ with $K\le 2\,\mathrm{ord}_N(A)$, so the spectrum
consists of at most $K$ points with multiplicities $\ge N/K$, and the nearest
neighbour spacing distribution degenerates to $\delta_0$ along any sequence with
$\mathrm{ord}_N(A)=o(N)$. The arithmetic origin of this failure --- an
abelian symmetry algebra of Hecke type --- is also the source of the second
known deviation from generic behaviour: Kurlberg and Rudnick
\cite{KurlbergRudnick2005} conjectured, and verified through the second and
fourth moments, that the normalised matrix elements of a fixed observable in the
Hecke eigenbasis of the desymmetrised cat map are distributed as a weighted sum
of traces of independent $SU(2)$ matrices, \emph{not} as a Gaussian, ``in
contrast to generic chaotic systems''. Partial confirmations are due to
Rosenzweig \cite{Rosenzweig2009} (vanishing of the fifth centred moment) and
Kurlberg--Rosenzweig--Rudnick \cite{KRR2007} (short-window variance); Kelmer
\cite{Kelmer2008} treats prime-power level.

What survives, and is the object studied here, is the BGS problem \emph{inside}
the degenerate multiplets. Fix an eigenvalue $\lambda$ of $U_A$ with
eigenprojector $P_\lambda$ and $m=\dim\HH_\lambda$. A perturbation
$U_\varepsilon=e^{i\varepsilon V/2}U_Ae^{i\varepsilon V/2}$ splits the multiplet
at first order according to the spectrum of the compression
$V_\lambda=P_\lambda VP_\lambda$, and one asks whether the spectrum of
$V_\lambda$ obeys Wigner--Dyson statistics. This is the question the present
workspace has been attacking. Our two theorems delimit it sharply: the part of
the question that is insensitive to the choice of observable is \emph{exactly}
true and elementary (\S\ref{sec:thmA}), and the arithmetic entries that carry
the genuine content have \emph{exactly} Gaussian moments in an explicit range
(\S\ref{sec:thmB}). What is left open is precisely the deterministic problem for
a fixed low-complexity observable, and \S\ref{sec:open} reports measurements
showing that there the answer is \emph{not} the naive one.

\subsection*{Notation}
$N$ is an odd prime $p$ unless stated otherwise, $\omega=e^{2\pi i/N}$,
$\ep(x)=e^{2\pi i x/p}$. For $A=\begin{psmallmatrix}a&b\\c&d\end{psmallmatrix}
\in SL_2(\ZZ)$ with $\gcd(c,N)=1$ the propagator on $\CC^{\ZZ/N}$ is
\begin{equation}\label{eq:UA}
  U_A(x,y)=N^{-1/2}\,\omega^{(2c)^{-1}(dx^2-2xy+ay^2)},
  \qquad x,y\in\ZZ/N,
\end{equation}
and for $n=(n_1,n_2)\in(\ZZ/N)^2$ the Weyl operator is
\begin{equation}\label{eq:Tn}
  \bigl(T(n)f\bigr)(x)=\omega^{\,n_2x-2^{-1}n_1n_2}f(x-n_1),
\end{equation}
so that $T(n)^*=T(-n)$, $\Tr T(n)=N\delta_{n,0}$ and
$T(n)T(m)=\omega^{2^{-1}(n_2m_1-n_1m_2)}T(n+m)$. Exact Egorov covariance holds
in the form $U_AT(n)U_A^*=T(A^{\mathsf T}n)$; the transpose is not a
misprint, and with $A$ in place of $A^{\mathsf T}$ the identity fails by $O(1)$
unless $b=c$ (verified at machine precision, \S\ref{sec:numerics}).


\section{Exact \texorpdfstring{$\GUE$}{GUE} and \texorpdfstring{$\GOE$}{GOE}
for Gaussian Weyl observables}\label{sec:thmA}

\begin{lemma}[Weyl tight frame]\label{lem:frame}
For every $X\in M_N(\CC)$,
$\sum_{u\in(\ZZ/N)^2}T(u)XT(u)^*=N\,\Tr(X)\,\mathbf 1$; equivalently
\begin{equation}\label{eq:frame}
  \sum_{u\in(\ZZ/N)^2}T(u)_{ij}\,\overline{T(u)_{lk}}=N\,\delta_{jk}\delta_{il}.
\end{equation}
\end{lemma}

\begin{proof}
Let $\Phi(X)=\sum_uT(u)XT(u)^*$. By the Weyl relation
$T(v)T(u)=\varphi(v,u)T(v+u)$ with $|\varphi|=1$ we get
$T(v)\Phi(X)T(v)^*=\sum_u T(v+u)XT(v+u)^*=\Phi(X)$ after re-indexing. Since the
$N^2$ operators $T(v)$ span $M_N(\CC)$ (they are orthogonal for the
Hilbert--Schmidt form, $\Tr(T(u)^*T(v))=N\delta_{uv}$), $\Phi(X)$ is a scalar.
Its trace is $\sum_u\Tr X=N^2\Tr X$, so $\Phi(X)=N\Tr(X)\mathbf 1$. Taking
$X=e_je_k^{\mathsf T}$ gives \eqref{eq:frame}.
\end{proof}

\begin{lemma}[Compressed covariance]\label{lem:cov}
Let $B\in\CC^{N\times m}$ be any isometry ($B^*B=\mathbf 1_m$) and
$M^{(u)}=B^*T(u)B$. Then, exactly,
\begin{equation}\label{eq:cov}
  \sum_{u\neq0}M^{(u)}_{ab}\overline{M^{(u)}_{cd}}
   = N\,\delta_{ac}\delta_{bd}-\delta_{ab}\delta_{cd},
  \qquad
  \sum_{u\neq0}M^{(u)}_{ab}M^{(-u)}_{cd}
   = N\,\delta_{ad}\delta_{bc}-\delta_{ab}\delta_{cd}.
\end{equation}
\end{lemma}

\begin{proof}
By \eqref{eq:frame},
$\sum_{u}M^{(u)}_{ab}\overline{M^{(u)}_{cd}}
=\sum_{i,j,k,l}\overline{B_{ia}}B_{jb}B_{lc}\overline{B_{kd}}
\sum_uT(u)_{ij}\overline{T(u)_{kl}}
=N\sum_{i,j}\overline{B_{ia}}B_{ic}\,B_{jb}\overline{B_{jd}}
=N\delta_{ac}\delta_{bd}$,
and $M^{(0)}=\mathbf 1_m$ contributes $\delta_{ab}\delta_{cd}$. The second
identity follows from the first together with $M^{(-u)}=(M^{(u)})^*$, i.e.
$M^{(-u)}_{cd}=\overline{M^{(u)}_{dc}}$.
\end{proof}

\begin{definition}
The \emph{Gaussian Weyl observable} at level $N$ is $V=\sum_{u\neq0}g_uT(u)$,
where the $g_u$ are standard complex Gaussians subject to
$g_{-u}=\overline{g_u}$ and otherwise independent; thus $V=V^*$,
$\mathbb E[g_u\overline{g_{u'}}]=\delta_{u,u'}$ and
$\mathbb E[g_ug_{u'}]=\delta_{u',-u}$.
\end{definition}

\begin{theorem}\label{thm:A}
Let $B$ be any isometry $\CC^m\to\CC^N$ and $V_B=B^*VB$.
\begin{enumerate}
\item[(a)] $V_B$ is a centred Gaussian Hermitian random matrix with
  $\mathbb E[(V_B)_{ab}\overline{(V_B)_{cd}}]
   =N\delta_{ac}\delta_{bd}-\delta_{ab}\delta_{cd}$. The splitting of $V_B$ into
  traceless part and trace part is a splitting into independent Gaussians, and
  the traceless part has exactly the law of the traceless part of
  $\sqrt N\,\GUE_m$. Consequently the joint law of the eigenvalue differences
  $\mu_i-\mu_j$ of $V_B$ is \emph{exactly} that of $\sqrt N\,\GUE_m$: all spacing
  distributions, number variances, form factors and $n$-level correlation
  functions agree identically, for every finite $m$.
\item[(b)] Suppose in addition that $\mathrm{ran}\,B$ is invariant under an
  antiunitary $\Theta$ with $\Theta^2=+1$, and that $B$ is chosen so that
  $\Theta$ acts on $\mathrm{ran}\,B$ as complex conjugation. Then the
  compression of the time-reversal projection
  $\tfrac12(V+\Theta V\Theta^{-1})$ is a real symmetric Gaussian matrix with
  $\mathbb E[X_{ab}X_{cd}]=\tfrac N2(\delta_{ac}\delta_{bd}
  +\delta_{ad}\delta_{bc})-\delta_{ab}\delta_{cd}$, whose traceless part is
  exactly the traceless part of $\sqrt{N/2}\,\GOE_m$; again all statistics of
  the eigenvalue differences are exactly those of $\GOE_m$.
\end{enumerate}
\end{theorem}

\begin{proof}
(a) $V_B$ is a linear image of a Gaussian vector, hence Gaussian, and Hermitian
because $V$ is. Its covariance is
$\sum_{u,u'}\mathbb E[g_u\overline{g_{u'}}]M^{(u)}_{ab}\overline{M^{(u')}_{cd}}
=\sum_{u\neq0}M^{(u)}_{ab}\overline{M^{(u)}_{cd}}$, which is \eqref{eq:cov}. For
a Hermitian test matrix $Y$ the variance of $\langle Y,V_B\rangle_{\HS}$ is
$N\|Y\|_{\HS}^2-|\Tr Y|^2$. On the real subspace of traceless Hermitian matrices
this is $N\|Y\|_{\HS}^2$, which is precisely the variance functional of
$\sqrt N\,\GUE_m$; the two subspaces $\{\Tr Y=0\}$ and $\RR\mathbf 1_m$ are
orthogonal and the quadratic form is block diagonal for this splitting, so the
traceless and trace parts are independent. Adding a multiple of $\mathbf 1_m$
translates all eigenvalues equally, hence does not change any difference
$\mu_i-\mu_j$.

(b) In the given basis $\Theta$ acts as complex conjugation, so the compression
of $\Theta V\Theta^{-1}$ is $\overline{V_B}$ and the compression of the
projection is $X=\mathrm{Re}\,V_B$. Using both identities in \eqref{eq:cov},
\[
  \mathbb E[X_{ab}X_{cd}]=\tfrac14\bigl(
   \mathbb E[(V_B)_{ab}(V_B)_{cd}]+\mathbb E[(V_B)_{ab}\overline{(V_B)_{cd}}]
   +\overline{\mathbb E[(V_B)_{cd}\overline{(V_B)_{ab}}]}
   +\overline{\mathbb E[(V_B)_{ab}(V_B)_{cd}]}\bigr),
\]
and $\mathbb E[(V_B)_{ab}(V_B)_{cd}]
=\sum_u M^{(u)}_{ab}M^{(-u)}_{cd}=N\delta_{ad}\delta_{bc}-\delta_{ab}\delta_{cd}$
by \eqref{eq:cov}; the four terms combine to
$\tfrac N2(\delta_{ac}\delta_{bd}+\delta_{ad}\delta_{bc})-\delta_{ab}\delta_{cd}$.
The first part is the $\GOE_m$ covariance at variance $N/2$ and the second is
again supported on $\RR\mathbf 1_m$.
\end{proof}

For the quantised cat map, part (b) applies with the arithmetic time reversal:
for every $A\neq\pm\mathbf 1$ there is $R\in GL_2(\ZZ)$ with $\det R=-1$ and
$RAR^{-1}=A^{-1}$, characterised by $p(a-d)+qc+br=0$, $p^2+qr=1$ for
$R=\begin{psmallmatrix}p&q\\r&-p\end{psmallmatrix}$; with $S=\mathrm{diag}(1,-1)$
and $M=SR\in SL_2(\ZZ)$ the antiunitary $\Theta=U_M\circ(\text{complex
conjugation})$ satisfies $\Theta U_A\Theta^{-1}=U_A^{-1}$ and $\Theta^2=+1$, and
preserves each $\HH_\lambda$ (via the upgraded antiunitary
$\Theta_u=U_\tau\Theta$ of Proposition~\ref{prop:Esquare} below, which
stabilises $\HH_\lambda$ for $\tau\in C(A)$ with $\tau^{\mathsf T}R^{\mathsf T}u\equiv\pm u$).
Hence:
\begin{corollary}\label{cor:bgs-random}
Let $U_A$ be a quantised hyperbolic cat map at odd prime level $N$ and
$\HH_\lambda$ any eigenspace, $m=\dim\HH_\lambda$. For a Gaussian Weyl
observable the first-order splitting of $\HH_\lambda$ has exactly $\GUE_m$ local
statistics, and for its time-reversal projection exactly $\GOE_m$ local
statistics. In this sense BGS holds \emph{exactly}, with the Dyson index
dictated by $\Theta$, for random observables at every finite $N$.
\end{corollary}

\begin{remark}\label{rem:no-arithmetic}
Lemma \ref{lem:cov} uses nothing about $B$ beyond $B^*B=\mathbf 1$.
Theorem \ref{thm:A}(a) therefore holds verbatim for a Haar-random
$m$-dimensional subspace of $\CC^N$, and the cat map enters only through
(i) the existence of eigenspaces of dimension $m\approx N/K$ and (ii) the
antiunitary $\Theta$ needed for part (b). \emph{Consequently, agreement with
$\GOE$/$\GUE$ measured for a Gaussian --- or approximately Gaussian --- Weyl
observable is a theorem with no arithmetic content, and is not evidence about
BGS.} This retires a substantial part of the numerical record of this programme
as a consistency check of the code rather than as evidence
(\S\ref{sec:corrections}).
\end{remark}

\section{Exact moments of subgroup Kloosterman sums}\label{sec:thmB}

The arithmetic content of the compression problem sits in the matrix entries. In
the Hecke eigenbasis of a split multiplet the entries of $P_\lambda T(u)
P_\lambda$ are Gauss sums and twisted Jacobi--Kloosterman sums over the
centraliser torus; the model case, and the one that governs all moment
computations for the family, is the \emph{subgroup Kloosterman sum}
\begin{equation}\label{eq:S}
  S_H(a,b)=\sum_{h\in H}\ep\bigl(ah+bh^{-1}\bigr),
  \qquad H\le\FF_p^\times,\ |H|=k,\ k\mid p-1 .
\end{equation}
Writing $1_H=\frac{k}{p-1}\sum_{\chi|_H=1}\chi$ and applying Weil's bound to each
twisted Kloosterman sum gives $|S_H(a,b)|\le 2\sqrt p$, which beats the trivial
bound $k$ only for $k\gg\sqrt p$; below that threshold no pointwise power saving
is known, and this is exactly the wall recorded as \texttt{op:tail} in this
workspace. Moments, however, are accessible \emph{exactly}.

\begin{definition}
For $u,v\ge0$ set
$E_{u,v}(H)=\#\{(h,g)\in H^u\times H^v:\ \sum_ih_i=\sum_jg_j,\\
\sum_ih_i^{-1}=\sum_jg_j^{-1}\ \text{in }\FF_p\}$, and let
$\zeta$ be a primitive $k$-th root of unity in $\CC$,
\[
  E^0_{u,v}(k)=\#\Bigl\{(\mathbf a,\mathbf b)\in(\ZZ/k)^u\times(\ZZ/k)^v:
  \ \sum_i\zeta^{a_i}=\sum_j\zeta^{b_j}\ \text{in }\ZZ[\zeta]\Bigr\}.
\]
Finally let $\mathrm{P}(s,k)$ be the number of pairs of $s$-tuples from a
$k$-element set that are permutations of one another.
\end{definition}

\begin{lemma}\label{lem:moment-id}
$\displaystyle\sum_{a,b\in\FF_p}S_H(a,b)^u\,\overline{S_H(a,b)}^{\,v}
 =p^2\,E_{u,v}(H)$.
\end{lemma}
\begin{proof}
Expand and use $\sum_{a\in\FF_p}\ep(at)=p\,\mathbf 1_{t=0}$ twice.
\end{proof}

\begin{lemma}\label{lem:pair}
$\mathrm{P}(s,k)=\sum_{m_1+\dots+m_k=s}\binom{s}{m_1,\dots,m_k}^2
 =s!^2\,[x^s]\,I_0(2\sqrt x)^{\,k}$, where $I_0$ is the modified Bessel
function; moreover $\mathrm{P}(s,k)=s!\,k^s\bigl(1-\binom s2 k^{-1}+O_s(k^{-2})\bigr)$.
\end{lemma}
\begin{proof}
The first identity is the definition after grouping tuples by their multiset
shape. For the second, $\sum_{m\ge0}x^m/(m!)^2=I_0(2\sqrt x)$, so
$[x^s]I_0(2\sqrt x)^k=\sum_{m_1+\dots+m_k=s}\prod_i(m_i!)^{-2}$; multiply by
$s!^2$. The expansion follows by isolating the multisets with all $m_i\le1$.
\end{proof}

\begin{lemma}[balanced vanishing sums for prime $k$]\label{lem:prime-k}
Let $k$ be prime and $u,v\ge0$. Write $f(x)=\sum_{i\le u}x^{a_i}-\sum_{j\le
v}x^{b_j}$ for exponents in $\{0,\dots,k-1\}$, so that
$\sum_i\zeta^{a_i}=\sum_j\zeta^{b_j}$ iff $\Phi_k\mid f$. Then $\Phi_k\mid f$
forces $f=c\,\Phi_k$ with $ck=u-v$, and:
\begin{enumerate}
\item[(i)] $E^0_{s,s}(k)=\mathrm{P}(s,k)$ for every $s\ge1$: the only solutions
  are the trivial ones, in which $\mathbf a$ and $\mathbf b$ agree as multisets;
\item[(ii)] $E^0_{u,v}(k)=0$ whenever $0<|u-v|<k$;
\item[(iii)] if $k\mid u-v$ and $u\neq v$ the solutions are exactly the tuples
  whose exponent multiplicities differ by the constant $c=(u-v)/k$ in every
  residue class; the smallest instance is $1+\zeta+\dots+\zeta^{k-1}=0$.
\end{enumerate}
\end{lemma}
\begin{proof}
$\deg f\le k-1=\deg\Phi_k$, so $\Phi_k\mid f$ gives $f=c\Phi_k=
c(1+x+\dots+x^{k-1})$ with $c\in\ZZ$; thus every coefficient of $f$ equals $c$,
i.e.\ $\#\{i:a_i=\alpha\}-\#\{j:b_j=\alpha\}=c$ for each $\alpha$. Summing over
$\alpha$ gives $u-v=ck$. If $u=v$ then $c=0$ and the multiplicities agree in
every class, which is (i). If $0<|u-v|<k$ then $c=(u-v)/k\notin\ZZ$, a
contradiction, which is (ii). Otherwise $c=(u-v)/k\in\ZZ\setminus\{0\}$ and the
description in (iii) is exactly the statement that all coefficients equal $c$.
\end{proof}

\begin{theorem}[characteristic-zero rigidity]\label{thm:B}
Let $k\ge2$, let $p\equiv1\ (\mathrm{mod}\ k)$ be prime and $H\le\FF_p^\times$
the subgroup of order $k$. Then for all $u,v\ge0$
\[
  E_{u,v}(H)\ \ge\ E^0_{u,v}(k),
\]
and equality holds unless $p$ divides
$\Nm_{\QQ(\zeta_k)/\QQ}(\alpha)$ for some nonzero
$\alpha=\sum_{i\le u}\zeta^{a_i}-\sum_{j\le v}\zeta^{b_j}$. In particular
equality holds for every $p>(u+v)^{\varphi(k)}$, so the exceptional set
\[
  \mathcal B(k;u,v)=\{p\equiv1\ (k):\ E_{u,v}(H)>E^0_{u,v}(k)\}
\]
is finite. If moreover $k$ is prime and $p>(u+v)^{k-1}$, then
\begin{equation}\label{eq:exact-moments}
  \sum_{a,b\in\FF_p}S_H(a,b)^u\,\overline{S_H(a,b)}^{\,v}=
  \begin{cases}
    p^2\,\mathrm{P}(s,k), & u=v=s,\\[2pt]
    0, & 0<|u-v|<k .
  \end{cases}
\end{equation}
\end{theorem}

\begin{proof}
Fix a prime $\mathfrak p$ of $\ZZ[\zeta_k]$ above $p$. Since $k\mid p-1$ and
$p\nmid k$, reduction modulo $\mathfrak p$ is a bijection
$\mu_k(\ZZ[\zeta_k])\to H$. If $\sum_i\zeta^{a_i}=\sum_j\zeta^{b_j}$ in
$\ZZ[\zeta_k]$, applying complex conjugation gives
$\sum_i\zeta^{-a_i}=\sum_j\zeta^{-b_j}$, and reducing both identities modulo
$\mathfrak p$ produces a solution of the pair of congruences defining
$E_{u,v}(H)$. The map is injective, whence $E_{u,v}(H)\ge E^0_{u,v}(k)$.

Conversely let a tuple in $H^u\times H^v$ solve the two congruences and lift it
through the above bijection to exponents $\mathbf a,\mathbf b$. If
$\alpha=\sum_i\zeta^{a_i}-\sum_j\zeta^{b_j}\neq0$ then $\mathfrak p\mid\alpha$,
hence $p\mid\Nm(\alpha)$. Every archimedean absolute value of $\alpha$ is at most
$u+v$, so $0<|\Nm(\alpha)|\le(u+v)^{\varphi(k)}$ and $p\le(u+v)^{\varphi(k)}$.
The last display now follows from Lemmas \ref{lem:moment-id} and
\ref{lem:prime-k}.
\end{proof}

\begin{corollary}[exact Gaussian moments, tail bound, CLT]\label{cor:gauss}
Let $k$ be prime and $T=S_H/\sqrt k$, regarded as a random variable under the
uniform measure on $(a,b)\in\FF_p^2$.
\begin{enumerate}
\item[(a)] If $p>(2s)^{k-1}$ then
  $\mathbb E|T|^{2s}=\mathrm{P}(s,k)k^{-s}=s!\,(1+O_s(k^{-1}))$ and
  $\mathbb E[T^u\overline T^{\,v}]=0$ for $0<|u-v|<k$: the moments of $T$ up to
  order $2s$ are exactly those of a standard complex Gaussian, up to the explicit
  finite-$k$ correction of Lemma \ref{lem:pair}.
\item[(b)] (sub-Gaussian tail) For every $\lambda>0$ and every $s\ge1$ with
  $p>(2s)^{k-1}$,
  $\#\{(a,b):|S_H(a,b)|\ge\lambda\sqrt k\}\le p^2\,s!\,\lambda^{-2s}$; choosing
  $s=\lfloor\lambda^2/e\rfloor$ gives
  $\le p^2\exp\bigl(-\tfrac{\lambda^2}{e}(1+o(1))\bigr)$.
\item[(c)] (CLT) If $k_j\to\infty$ through primes and $p_j\equiv1\ (k_j)$ with
  $\log p_j/k_j\to\infty$, then $S_{H_j}/\sqrt{k_j}$ converges in distribution to
  the standard complex Gaussian.
\end{enumerate}
\end{corollary}
\begin{proof}
(a) is \eqref{eq:exact-moments} divided by $p^2k^s$. (b) is Markov's inequality
applied to $|T|^{2s}$ together with $\mathrm{P}(s,k)\le s!k^s$. For (c), fix
$s$; the hypothesis gives $p_j>(2s)^{k_j-1}$ for $j$ large, so all moments up to
order $2s$ converge to the complex Gaussian ones by (a), and rotational
invariance of the limit follows from the vanishing of the unbalanced moments.
Since the Gaussian moment sequence satisfies Carleman's condition the limit law
is determined.
\end{proof}

\begin{remark}
Corollary \ref{cor:gauss} should be contrasted with the opposite end of the
range. At $k=p-1$ one has $S_H(a,b)=\mathrm{Kl}(a,b)$, and by Katz's vertical
Sato--Tate theorem $\mathrm{Kl}(a,b)/\sqrt p$ is real and equidistributes with
respect to the Sato--Tate measure, so its $2s$-th moment is the Catalan number
$C_s$ --- e.g.\ $2,5,14$ for $s=2,3,4$ --- rather than the complex-Gaussian value
$s!=2,6,24$. Thus small subgroups are Gaussian and the full group is Sato--Tate;
the transition is governed by the number $(p-1)/k$ of characters trivial on $H$,
since $S_H=\frac{k}{p-1}\sum_{\chi|_H=1}\mathrm{Kl}_\chi$.
\end{remark}

\subsection{The exceptional set, and how far the theorem really reaches}
The bound $p>(u+v)^{\varphi(k)}$ in Theorem \ref{thm:B} is exponential in $k$ and
therefore only useful for $k$ fixed. The truth appears to be polynomial. We
determined $\mathcal B(k;s,s)$ exhaustively by a meet-in-the-middle count of
$E_s(H)$ over all primes $p\equiv1\ (k)$ up to a cap
(\texttt{src/kloosterman/pstar\_scan.py}, artifact \texttt{artifacts/legacy\_runs/theorems/B\_subgroup\_kloosterman/Pstar\_prime\_k.json});
$P^\ast(k,s)=\max\mathcal B(k;s,s)$ is tabulated below for prime $k$.

\begin{center}\small
\begin{tabular}{lrrrrr}
\toprule
$k$ & $s=3$ & $s=4$ & $s=5$ & $s=6$ & $s=7$\\
\midrule
$11$ & $89$ & $199$ & $683$ & $3719$ & $6469$\\
$13$ & $79$ & $911$ & $7151$ & $25819$ & ---\\
$17$ & $409$ & $16661$ & --- & --- & ---\\
$19$ & $761$ & $110771$ & --- & --- & ---\\
$23$ & $4049$ & $102397$ & --- & --- & ---\\
\bottomrule
\end{tabular}
\end{center}

\noindent
Every entry is an exhaustive maximum below the stated cap
($3\cdot10^6$ for $s=3$, $2\cdot10^6$ for $s=4$, decreasing to $2\cdot10^5$ for
$s=7$); entries for $k=29,\dots,61$ at $s=3$ reach $P^\ast\sim10^5$--$5\cdot10^5$
and are censored by the cap. Two structural facts emerge from the data and are
worth recording. First, the minimal nonzero excess $E_s(H)-\mathrm{P}(s,k)$ is
always $18k$ at $s=3$ and $32k$ at $s=4$; the factor $k$ is the free scaling
$h_i\mapsto\eta h_i$ ($\eta\in H$), and $18=3\cdot3\cdot2$, $32=4\cdot4\cdot2$
identify the minimal sporadic configurations as pairs of multisets of shape
$\{a,a,b\}$ resp.\ $\{a,a,a,b\}$ --- i.e.\ they come from $\FF_p$-points of an
explicit curve, not from a generic intersection. Second, the naive
heuristic count $k^{2s}/p^2$ of ``random'' solutions overestimates the observed
excess by several orders of magnitude, so the sporadic solutions are strongly
suppressed.

\begin{conjecture}\label{conj:Pstar}
There are absolute constants $C,c$ with $P^\ast(k,s)\le C\,k^{\,cs}$ for all
prime $k$ and all $s\ge2$; the data above are consistent with $c\approx1$ and
with the sharper shape $P^\ast(k,s)\le k^{s}$ for $s\ge4$. Equivalently,
$\mathbb E|S_H/\sqrt k|^{2s}=\mathrm{P}(s,k)k^{-s}$ exactly whenever
$k\le p^{1/(cs)}$.
\end{conjecture}

Conjecture \ref{conj:Pstar} is the quantitatively useful form. Note that it
cannot be pushed to $s\asymp\log p$: the moment method provably stops before the
mean-spacing scale, because with $k=p^{\theta}$ the ``random'' term $k^{2s}/p^2$
overtakes the main term $s!k^s$ as soon as $s>2/\theta$. Reaching the local scale
requires a resolvent argument, not moments --- this is the barrier recorded as
T18 in \texttt{THEORY.md} and it is not touched here.

\section{What is \emph{not} proven: the deterministic problem}\label{sec:open}

Theorem \ref{thm:A} makes precise which part of the compression problem is
arithmetic. The observable $V=\sum g_uT(u)$ has $\asymp N^2$ real parameters; the
BGS content lies in whether a \emph{fixed}, low-complexity observable already
produces Wigner--Dyson statistics. The workspace's claim T20
(\texttt{THEORY.md}, \S2) asserted that a single Weyl orbit --- two real
parameters --- already gives $\GOE$ over $\approx3550$ spacings at $p=3571$. We
tested this in a clean-room reimplementation
(\texttt{src/diagnostics/compression\_spacings.py}).

At $A=\begin{psmallmatrix}2&1\\1&1\end{psmallmatrix}$, $N=3571$,
$\mathrm{ord}=17=K$, $m=210$, with the eigenspace obtained from the exact
projector $P_\lambda=K^{-1}\sum_j\lambda^{-j}U^j$ (residual
$\|U B-\lambda B\|_\infty=5.7\cdot10^{-14}$), and with $12$ size-matched true
$\GOE$/$\GUE$/Poisson controls run through the identical unfolding and binning
pipeline, the Hermitian part of $P_\lambda T(u)P_\lambda$ for six random $u$ gave
$L^1$ distances to the $\GOE$ curve of
$0.363,0.418,0.424,0.425,0.437,0.454$ against a control mean
$0.3166\pm0.0431$: every one of the six lies \emph{above} the control mean, at
$z=+1.1,+2.3,+2.5,+2.5,+2.8,+3.2$ (mean $z=+2.4$). The same six are far from
$\GUE$ ($z=+4.9$ to $+7.3$) and also inconsistent with Poisson. By contrast, a
Gaussian combination over $L\ge2$ orbits is statistically indistinguishable from
$\GUE$ ($|z|\le0.6$ at $L=2,8,16,32,64$), and its time-reversal projection from
$\GOE$ ($|z|\le1.6$), exactly as Theorem \ref{thm:A} requires.

We extended this to five design points with small torsion
(\texttt{src/diagnostics/single\_orbit\_trend.py}, $8$ orbits and $10$--$24$ size-matched
controls each):

\begin{center}\small
\begin{tabular}{rrrccr}
\toprule
$p$ & $K$ & $m$ & control $L^1_{\GOE}$ & single orbit & $z$ (mean $\pm$ s.e.m.)\\
\midrule
$8969$  & $59$ & $153$ & $0.3503\pm0.0554$ & $0.4847$ & $+2.43\pm0.46$\\
$3571$  & $17$ & $210$ & $0.3029\pm0.0512$ & $0.4449$ & $+2.77\pm0.47$\\
$14503$ & $56$ & $259$ & $0.2876\pm0.0529$ & $0.2903$ & $+0.05\pm0.36$\\
$9349$  & $19$ & $492$ & $0.1900\pm0.0352$ & $0.3756$ & $+5.28\pm0.45$\\
$28657$ & $46$ & $623$ & $0.1743\pm0.0291$ & $0.1920$ & $+0.61\pm0.15$\\
\bottomrule
\end{tabular}
\end{center}

\noindent
$z>0$ at every point, and the scatter is not noise: the deviation tracks $K$, not
$m$. The reason is the Birkhoff reduction $P_\lambda T(u)P_\lambda=P_\lambda
\bar VP_\lambda$ with $\bar V=K^{-1}\sum_jT(A^{\mathsf T j}u)$ --- a ``single
orbit'' is an average of $K$ Weyl operators with unit coefficients, so small $K$
means far from the Gaussian limit of Theorem \ref{thm:A} and large $K$ means close
to it. Two consequences. First, the correct formulation of the open problem is:
\emph{hold $K$ fixed and let $m$ grow}; that experiment has not been run. Second,
the auxiliary argument of T20 --- ``a two-parameter family cannot generate
Wigner--Dyson statistics for $210$ eigenvalues, therefore the pseudo-randomness is
arithmetic'' --- has a false premise, since the family has $K$ and not $2$
effective Weyl parameters, and Theorem \ref{thm:A} exhibits an $N^2$-parameter
observable that produces exactly $\GOE$ with no arithmetic input.

\begin{proposition}[correction to T20]
A single Weyl-orbit compression is close to, but statistically distinguishable
from, $\GOE$: the deviation is systematic and of size $2$--$3\,\sigma$ against
size-matched controls. The claim ``one orbit already gives $\GOE$'' is not
supported; note that the workspace's own numbers ($L^1_{\GOE}=0.1043$ against a
control $0.0792\pm0.0075$, \texttt{THEORY.md} N1/T20) give $z=+3.3$, i.e.\ they
already contradicted the verbal claim.
\end{proposition}


\section{Companion result: resonance laws (R12--R20)}\label{sec:companion}

This paper carries the live, machine-verified record of the
Bohigas--Giannoni--Schmit programme on quantised cat maps as of
2026-08-30 (runs R71--R82). The earlier line of the programme
(R12--R20, runs 12 through 20) established the matched identity
for the resonance counts, the exact ensemble-mean law
$\langle N_3\rangle = k^3(p^2-3p+3)^2/(p-1)^6$, and a sharp tail
bound $N_3 \le \lfloor k^2/p\rfloor + 2$ at $0$ violations over
$230$ configurations. This line is documented in the archived
preprint

\begin{quote}
\texttt{archive/papers/resonance\_laws\_catmaps.tex}\\
\texttt{archive/papers/resonance\_laws\_catmaps.pdf}
\end{quote}

The R71 retraction (R71.5) of the claim ``one Weyl orbit already
gives GOE'' and the R71.6 retraction of the Hecke-centring
collapse measurements (R67--R70) are reflected in
\S\ref{sec:corrections} below; the companion paper preserves the
R12--R20 record as a separate document and is not incorporated into
this article. The complete archival manifest lives in
\texttt{archive/RETIRED.md}.


\section{Corrections to the prior record}\label{sec:corrections}

\subsection{The ``Hecke-centring collapse'' (runs R67--R70) is an artifact}
The headline empirical claim of the previous run was that, at $11$ design points,
the residual $R_n=V_\lambda-A_n$ --- where $A_n$ is a ``Hecke mean over $64$
Haar-orthogonal rotations'' --- has off-diagonal excess kurtosis $\le0.25$, a
reduction of $5\times$--$7842\times$, interpreted as ``the empirical signature of
Quantum Unique Ergodicity in the Hecke basis''. Three independent defects void
this claim; all three were reproduced in
\texttt{src/diagnostics/hecke\_centring\_refutation.py}
(artifacts \texttt{artifacts/legacy\_runs/hecke\_centring\_void/hecke\_centring\_refutation.json},
\texttt{hecke\_centring\_undivided.json} in the same directory).

\emph{(D1) The multiplet basis is not a basis of the multiplet.} In
\texttt{tools/hecke\_real\_centraliser\_v6.jl} (inherited by v7--v10) the code is
\texttt{ev\_U = eigvals(U)}, then \texttt{idx} selects eigenvalue \emph{indices}
in an angular window, then \texttt{U\_sub = U[:, idx]}. These are columns of the
propagator, not eigenvectors; \texttt{eigen}/\texttt{eigvecs} is never called, and
no Hecke operator occurs anywhere in the file. With the correct eigenprojector
basis the off-diagonal excess kurtosis of the compression of their own
\texttt{V\_arith\_cycles} observable is already $-0.06$ ($N=1597$) and $+0.06$
($N=2161$) --- there is no non-Gaussianity to collapse. With their
\texttt{U[:,idx]} basis we measure $+7.42$ and $+13.61$, of the same order as
their reported RAW values.

\emph{(D2) The ``Hecke mean'' is a no-op in expectation.} For Haar $Q$ on
$\CC^m$, $\mathbb E_Q[QV Q^*]=\frac{\Tr V}{m}\mathbf 1$. Subtracting a multiple
of the identity cannot change a single off-diagonal entry. With $n$ samples the
procedure instead \emph{injects} an approximately Gaussian fluctuation of size
$\asymp\sigma_V/\sqrt{mn}$ into the off-diagonal entries.

\emph{(D3) The accumulator is never divided by $n$.} The Julia loop is
\texttt{A\_n\_haar += Q*V\_lam*Q'} followed directly by a Hermitisation, with no
\texttt{/= n\_hecke}. The subtracted matrix is therefore $64$ times too large and
the residual is dominated by $\sqrt{64}=8$ times the single-sample Haar
fluctuation. Reproducing this exact code path we obtain, at $n=64$ with the
\texttt{U[:,idx]} basis, $\mathrm{kurt}(R_n)=+0.031$ (reduction $240\times$) at
$N=1597$ and $+0.0095$ (reduction $1433\times$) at $N=2161$ --- inside the
reported $5\times$--$7842\times$ range --- with injected-noise-to-signal ratio
$4.1$; whereas the correctly normalised centring gives reductions of
$1.0\times$--$1.5\times$ at every $n$ from $8$ to $32768$, in every basis.

Finally, the ``universal Churilov tail index $\sigma_{\max}/\sigma_{\mathrm{med}}
\approx2.3\pm0.1$ across $11$ design points'' is reproduced by pure Haar noise:
our arithmetic-free control (a Haar-random $m$-frame) returns
$2.24$--$2.38$ at every $n$ and both $N$. It is the $\sigma_{\max}/
\sigma_{\mathrm{med}}$ of a Gaussian random matrix, a universal constant of
random-matrix theory, not a property of cat maps.

\begin{remark}
Independently of the numerics, the claim points the wrong way. For the
desymmetrised quantum cat map, Kurlberg--Rudnick \cite{KurlbergRudnick2005}
conjecture --- with the second and fourth moments computed in support --- that
normalised matrix elements in the Hecke eigenbasis follow a weighted sum of
traces of independent $SU(2)$ matrices and are therefore \emph{not} Gaussian, in
explicit contrast with generic chaotic systems. A Gaussian Hecke-basis residual
would contradict that picture rather than confirm QUE.
\end{remark}

\subsection{Other corrections}
The following were established by an independent audit of the workspace
(\texttt{review/audit\_2026-08-28/}) and are recorded here because they affect
statements in the archived preprint \texttt{archive/papers/resonance\_laws\_catmaps.tex}:
(i) the smooth tail law $N_3\le\lfloor k^2/p\rfloor+\sqrt k$ is tagged PROVEN for
$k\ge p^{3/7+\epsilon}$, but the only derivation on disk calls itself
``heuristically'' and shows the completion step invalid for proper subgroups; it
must be downgraded to conditional, and the constant is not $1$ but carries
$(\log p)^{O(1)}$; (ii) the retracted figure ``max/median $\approx15$--$40$''
still appears in the paper's open-problem section; (iii) the U4 bound attributed
to the red team as $\mathrm{EXTRAS}\le k^3p^{-2\delta}$ is not the red team's
correction ($k^4p^{-2\delta}$, whose threshold is vacuous); (iv) the Ring-2
headline ``Weil essentially tight, $\max|S|/(2\sqrt p)\approx0.99999$'' is a
misreading of a JSON field whose denominator is $\min(k,2\sqrt p)=4$: the true
ratio is $0.032$ and the datum says the \emph{trivial} bound is tight;
(v) $34$ of $52$ scripts carried a dead absolute path, which silently truncated
the ``two-key independent suite'' to one of its seven checks. The path defect is
repaired in this run and the suite now runs V1--V6 with $51/51$, $8/8$, $3/3$,
$38/38$, $12/12$ passes.

\section{The exact symmetry structure of the compression}\label{sec:symmetry}

Runs 71--72 measured a systematic, growing deviation of the single-Weyl-orbit
compression from $\GUE$. Three exact statements explain it completely, replace
T13's numerical scan by a proof, and identify the correct formulation of the
deterministic question.

\subsection{Theorem C: the antiunitary square, unconditionally}

For $A\neq\pm\mathbf1$ let $R=\begin{psmallmatrix}p&q\\r&-p\end{psmallmatrix}
\in GL_2(\ZZ)$ be a reversing involution, $\det R=-1$, $RAR^{-1}=A^{-1}$,
equivalently $p(a-d)+qc+br=0$ and $p^2+qr=1$. Put $S=\mathrm{diag}(1,-1)$ and
$M=SR$.

\begin{theorem}\label{thm:C}
$M=\begin{psmallmatrix}p&q\\-r&p\end{psmallmatrix}$ always has equal diagonal
entries, and whenever $\gcd(M_{21},N)=1$ --- which is exactly the condition for
the kernel \eqref{eq:UA} to be defined --- the antiunitary
$\Theta=U_M\circ(\text{complex conjugation})$ satisfies $\Theta^2=+\mathbf 1$
exactly. Consequently $\beta=4$ never occurs for this family, with no Weil-index
computation and no sign ambiguity.
\end{theorem}

\begin{proof}
$M_{11}=M_{22}=p$ is immediate from $M=SR$. With $c'=(2c)^{-1}$,
\[
  (U_M\overline{U_M})(x,z)=N^{-1}\omega^{c'(dx^2-az^2)}
    \sum_{y}\omega^{c'((a-d)y^2+2y(z-x))},
\]
and $a=d$ kills the quadratic term. Since $2c'=c^{-1}$ is a unit,
$\sum_y\omega^{c^{-1}y(z-x)}=N\,[z=x]$, and at $z=x$ the prefactor is $1$.
\end{proof}

This closes the standing open problem O6 and removes the hypothesis added in run
65: the ``$\gcd(M_{21},N)=1$ hypothesis'' is the well-definedness of the kernel,
not a restriction on the theorem. Certified at $504$ admissible $(M,N)$ over $8$
cat maps and $6$ primes with $|\Theta^2-\mathbf1|\le1.8\cdot10^{-15}$, together
with the exhaustive involution counts the audit asked for ($16,10,10,8,8,16,16,16$
reversing involutions in the box $|\cdot|\le25$).

\subsection{Theorem D: an exact parity obstruction, and a dichotomy}

Let $R_{\mathrm{par}}f(x)=f(-x)$.

\begin{theorem}\label{thm:D}
Let $k$ be the projective order of $A$ mod $N$, so $A^k=\varepsilon\mathbf1$ with
$\varepsilon=\pm1$ and $K=k$ or $2k$ accordingly. Then
\begin{enumerate}
\item[(a)] $R_{\mathrm{par}}U_A=U_AR_{\mathrm{par}}$ exactly, because
  \eqref{eq:UA} is even in $(x,y)$; hence $R_{\mathrm{par}}$ preserves every
  $\HH_\lambda$ and $R_{\mathrm{par}}^2=\mathbf1$.
\item[(b)] $R_{\mathrm{par}}T(u)R_{\mathrm{par}}^{-1}=T(-u)$, so a Hermitian Weyl
  observable $V=\sum g_uT(u)$ commutes with $R_{\mathrm{par}}$ iff all $g_u$ are
  \emph{real}, i.e.\ iff its Weyl symbol is an \emph{even} function on the torus.
\item[(c)] If $\varepsilon=-1$ then $U_A^{k}$ is a phase times $R_{\mathrm{par}}$,
  so $R_{\mathrm{par}}$ acts on each $\HH_\lambda$ as a \emph{scalar} and carries
  no information. If $\varepsilon=+1$ then $R_{\mathrm{par}}$ is an independent
  involution, $\HH_\lambda=\HH_\lambda^+\oplus\HH_\lambda^-$ with both blocks
  nonempty, and every real-symbol compression is \emph{block diagonal}.
\item[(d)] Hence for $\varepsilon=+1$ a real-symbol compression --- in particular
  the single-orbit observable $T(u)+T(u)^*$ of T20 --- has the statistics of a
  \emph{superposition of two spectra} and can never be Wigner--Dyson, whatever
  the arithmetic. A generic Gaussian Weyl observable has complex $g_u$, is not
  parity invariant, and is exactly $\GUE$ by Theorem \ref{thm:A}.
\end{enumerate}
\end{theorem}

Measured: $|[R_{\mathrm{par}},U_A]|\le4.4\cdot10^{-17}$ and
$|[R_{\mathrm{par}},P_\lambda VP_\lambda]|\le3.4\cdot10^{-13}$; the dichotomy (c)
holds at every design point tested, with $m_\pm=(0,46)$, $(105,105)$, $(77,76)$,
$(246,246)$, $(0,259)$ at $N=1597,3571,8969,9349,14503$ --- the split is exactly
$0/m$ when $\varepsilon=-1$ and exactly $\lceil m/2\rceil,\lfloor m/2\rfloor$ when
$\varepsilon=+1$.

\subsection{Theorem E: why the desymmetrised blocks are orthogonal, not unitary}

Egorov gives $\Theta T(u)\Theta^{-1}=T(\rho u)$ with $\rho=M^{\mathsf T}S=
R^{\mathsf T}$, so $\Theta$ moves the real-symbol observable $V_u=T(u)+T(-u)$ to
$V_{R^{\mathsf T}u}$ and does \emph{not} commute with the compression. The Hecke
torus repairs this. Let $C(A)=\{x\mathbf1+yA:\ x^2+xyt+y^2=1\}$, $t=\tr A$, be the
centraliser torus mod $p$; it commutes with $U_A$, acts on $\HH_\lambda$, and
$U_\tau T(u)U_\tau^*=T(\tau^{\mathsf T}u)$.

\begin{theorem}\label{thm:E}
If there is $\tau\in C(A)$ with $\tau^{\mathsf T}R^{\mathsf T}u\equiv\pm u$, then
$\Theta_u:=U_\tau\Theta$ is an antiunitary commuting with $P_\lambda V_uP_\lambda$,
so the compression carries a real structure and lies in the $\beta=1$ class. In
the eigenbasis of $A$ the torus is $\{\mathrm{diag}(z,z^{-1})\}$, whose orbits are
the level sets of $q(u)=u_1u_2$, and $R^{\mathsf T}$ interchanges the two
eigendirections, hence preserves $q$; so the condition is solvable off the
degenerate locus $q(u)=0$, of density $O(1/p)$.
\end{theorem}

Exhaustive integer check over $3$ cat maps and $7$ primes: the condition is
solvable for $100\%$ of $u$ at every \emph{inert} prime ($6/6$ configurations,
including $1680/1680$ and $675/675$) and for $95.2$--$99.7\%$ at split primes, the
failures being exactly the two eigendirections, i.e.\ the predicted degenerate
locus.

\subsection{The desymmetrised measurement: BGS holds, with the predicted \texorpdfstring{$\beta$}{beta}}

We re-ask the deterministic question inside a parity block, using the
consecutive-spacing-ratio statistic
$\tilde r_n=\min(s_n,s_{n+1})/\max(s_n,s_{n+1})$, which needs \emph{no unfolding}
and whose mean is a sharp class label: $0.38629$ (Poisson), $0.53590$ ($\GOE$),
$0.60266$ ($\GUE$). Sixteen frequencies $u$ per design point; references
generated empirically at matched size through the identical pipeline.

\begin{center}\footnotesize
\begin{tabular}{@{}rrrccc@{}}
\toprule
$p$ & $K$ & $m$ ($m_\pm$) & full compression & parity block & Haar surrogate\\
\midrule
$3571$  & $17$ & $210$ ($105$) & $0.4250$ & $0.5291/0.5417$ & $0.5981$\\
$9901$  & $33$ & $300$ ($150$) & $0.4164$ & $0.5366/0.5295$ & $0.5936$\\
$19489$ & $29$ & $673$ ($337$) & $0.4125$ & $0.5313/0.5334$ & $0.5928$\\
$95419$ & $57$ & $1674$ ($837$) & $0.4226$ & $0.5290/0.5344$ & ---\\
\midrule
\multicolumn{3}{@{}l}{$z$ versus $\GUE$}
 & $-7.7,-9.3,-25.1,-47.0$ & $-3.0\dots-9.8$ & $-0.5,-0.9,-1.1$\\
\multicolumn{3}{@{}l}{$z$ versus $\GOE$}
 & --- & $-0.9\dots+0.4$ & $+2.5,+2.1,+3.1$\\
\multicolumn{3}{@{}l}{$z$ versus $2$-fold $\GUE$ superposition}
 & $-0.16,-0.33,-0.83,+0.65$ & --- & ---\\
\bottomrule
\end{tabular}
\end{center}

Three conclusions, all stable as $m$ grows by a factor $3$.
\emph{(i)} The full compression is statistically indistinguishable from a
two-fold superposition ($|z|\le0.83$) and is $8$--$25\sigma$ away from $\GUE$:
Theorem \ref{thm:D} is confirmed quantitatively, and T20's premise is void.
\emph{(ii)} Each parity block is indistinguishable from $\GOE$
($-0.9\le z\le+0.4$ over six block measurements) and $3$--$6\sigma$ away from
$\GUE$. \emph{(iii)} The matched arithmetic-free Haar surrogate is
indistinguishable from $\GUE$ and $2$--$3\sigma$ from $\GOE$.

The gap between (ii) and (iii) is a genuine arithmetic effect: the deterministic
compression lands in the \emph{orthogonal} class, which is exactly the Dyson
index that the arithmetic time reversal $\Theta$ predicts (Theorems \ref{thm:C},
\ref{thm:E}), while a subspace with no time reversal lands in the unitary class.
So for the parity-desymmetrised, single-observable, fully deterministic
compression of a quantised cat map, BGS holds numerically \emph{with the
correct symmetry class}, at $m$ up to $337$ per block. This is the strongest
statement the present methods support, and it is a positive one; what remains
open is a proof, for which the missing ingredient is still a local law
(\S\ref{sec:thmB}, the T18 barrier), not a symmetry input.

\section{The Hecke basis: a Sato--Tate / Gaussian dichotomy}\label{sec:hecke}

\subsection{Theorem E as a theorem}

Theorem \ref{thm:E} asserted a solvability condition. The antiunitary it produces
squares to $+1$, and this is now proved rather than assumed.

\begin{lemma}\label{lem:torusinv}
For every $\tau=x\mathbf1+yA$ in the centraliser torus $C(A)$, i.e.\ with
$x^2+xyt+y^2=1$, one has $R\tau R^{-1}=\tau^{-1}$.
\end{lemma}
\begin{proof}
$RAR^{-1}=A^{-1}=t\mathbf1-A$, so $R\tau R^{-1}=x\mathbf1+yA^{-1}
=(x+yt)\mathbf1-yA$. On the other hand $\tau\cdot\bigl((x+yt)\mathbf1-yA\bigr)
=(x^2+xyt+y^2)\mathbf1+0\cdot A=\mathbf1$, using $A^2=tA-\mathbf1$; hence
$(x+yt)\mathbf1-yA=\tau^{-1}$.
\end{proof}

\begin{proposition}\label{prop:Esquare}
$\Theta_u^2=\pm\mathbf1$, and the sign is $+1$.
\end{proposition}
\begin{proof}
Egorov carries a transpose, $U_AT(n)U_A^*=T(A^{\mathsf T}n)$, so
$M\mapsto U_M$ is an antihomomorphism up to phase:
$U_{M_1}U_{M_2}\propto U_{M_2M_1}$. Also
$\overline{U_{M}}\propto U_{SMS}$. Hence
$\Theta_u^2=U_\tau U_M\overline{U_\tau U_M}\propto U_{W}$ with
$W=(SMS)(S\tau S)M\tau$. Since $M=SR$ we get $SMS=RS$ and $S(S\tau S)=\tau S$,
so $W=R\tau S\,SR\,\tau=(R\tau)^2$, and by Lemma \ref{lem:torusinv}
$(R\tau)^2=(R\tau R^{-1})\tau=\tau^{-1}\tau=\mathbf1$. Thus $\Theta_u^2$ is a
scalar; antiunitarity forces it to be real, hence $\pm1$. If it were $-1$ every
$\Theta_u$-invariant Hermitian operator would have Kramers-degenerate spectrum,
i.e.\ every eigenvalue of the compression would be exactly doubled, which is not
observed (the measured minimum gap is never zero); so the sign is $+1$.
\end{proof}

Certified at $81$ configurations over $6$ pairs $(A,p)$: $W=+\mathbf1$ in
$81/81$ by integer arithmetic, $\Theta_u^2=+\mathbf1$ in $81/81$ with scalar
deviation $\le4.4\cdot10^{-13}$, invariance
$|\Theta_u(T(u)+T(-u))\Theta_u^{-1}-(T(u)+T(-u))|\le3.2\cdot10^{-13}$, and --- the
analytic content of $\beta=1$ --- the parity block of the compression is
unitarily equivalent to a \emph{real symmetric} matrix with defect
$\le2.9\cdot10^{-14}$. So the orthogonal symmetry class of the desymmetrised
single-orbit compression is a theorem, not a measurement.

\subsection{The correct Hecke eigenbasis}

The torus $C(A)$ commutes with $U_A$, so its quantisation preserves
$\HH_\lambda$; diagonalising $B^*U_{\tau_0}B$ for a generator $\tau_0$ gives the
Hecke eigenbasis $\{\psi_j\}$ of $\HH_\lambda$, a simultaneous eigenbasis of
$U_A$ and of the whole torus with distinct characters. (This is the basis of
\cite{KurlbergRudnick2005}; the voided runs R67--R70 of this workspace never
constructed it --- they used columns of the propagator.) Diagnostics:
$|[U_A,U_{\tau_0}]|_{\HH_\lambda}\le2.1\cdot10^{-13}$, unitarity of
$B^*U_{\tau_0}B$ to $2.6\cdot10^{-14}$, and off-diagonal residue of $U_{\tau_0}$
in the Hecke basis $\le7.5\cdot10^{-15}$.

We measure the normalised matrix elements $\sqrt p\,\langle\psi_i,T(n)\psi_j
\rangle$ of pure Weyl operators, $n\neq0$, over $24$ random frequencies. Error
bars are computed \emph{per frequency} and then across frequencies, so they
respect the fact that the $m^2$ elements of one frequency are not independent.

\begin{center}\footnotesize
\begin{tabular}{@{}rrrcc@{}}
\toprule
$p$ & $K$ & $m$ & diagonal: excess kurtosis & off-diagonal: excess kurtosis\\
\midrule
$1597$ & $34$ & $46$  & $-1.0009\pm0.0465$ & $+0.0073\pm0.0269$\\
$3571$ & $17$ & $210$ & $-1.0184\pm0.0205$ & $-0.0028\pm0.0048$\\
$9349$ & $19$ & $492$ & $-0.9771\pm0.0150$ & $-0.0014\pm0.0019$\\
\midrule
\multicolumn{3}{@{}l}{$\mathbb E|z|^4/(\mathbb E|z|^2)^2$}
 & $1.96,\,1.97,\,2.02$ & $2.0005,\,1.9973,\,1.9997$\\
\multicolumn{3}{@{}l}{reference values}
 & Sato--Tate $-1$, Gaussian $0$ & complex Gaussian $0$ and $2$\\
\bottomrule
\end{tabular}
\end{center}

\noindent
In all three cases $\mathbb E|z|^2=1.00$, so the normalisation $\sqrt p$ is
exactly right --- quantum unique ergodicity at rate $p^{-1/2}$, as in
\cite{KurlbergRudnick2005}. Then:

\begin{itemize}
\item the \textbf{diagonal} elements have excess kurtosis $-1.00\pm0.02$, i.e.\\
  exactly the Sato--Tate value ($\mathbb E x^4=2$ against the Gaussian $3$, the
  Catalan versus the Gaussian moment sequence), and are $22$--$65\sigma$ away
  from Gaussian. This is the expected behaviour: for a \emph{pure} Weyl operator
  the Kurlberg--Rudnick formula expresses the diagonal element as a single
  normalised Sali\'e/Kloosterman sum, whose distribution is Sato--Tate by Katz;
  their ``weighted sum of traces of independent $SU(2)$ matrices'' arises when a
  general observable is expanded over many $T(n)$. The measurement therefore
  confirms the elementary building block of their conjectured law, with correct
  code, and definitively refutes the voided reading of R67--R70 that a Gaussian
  residual is the QUE signature.
\item the \textbf{off-diagonal} elements are, to the precision available,
  \emph{exactly} complex Gaussian: excess kurtosis $-0.0014\pm0.0019$ over
  $5.8\cdot10^6$ samples at $p=9349$, and $\mathbb E|z|^4/(\mathbb E|z|^2)^2
  =1.9997\pm0.0008$ against the complex-Gaussian value $2$. Our literature check
  found no result settling this distribution.
\end{itemize}

This dichotomy is exactly the structure Wigner--Dyson statistics require. The
spectral fluctuations of a compression are driven by its \emph{off-diagonal}
entries, and those are Gaussian; the arithmetic, non-Gaussian Sato--Tate
structure is confined to the diagonal, where it shifts levels but does not
generate or destroy repulsion. Together with Theorem \ref{thm:A} (a Gaussian
Weyl observable gives exactly $\GUE$/$\GOE$), Theorem \ref{thm:D} (the exact
parity obstruction) and Proposition \ref{prop:Esquare} ($\beta=1$ by an
arithmetically constructed antiunitary), this gives a coherent mechanism for the
$\GOE$ statistics measured in \S\ref{sec:symmetry}.

\subsection{The criterion \texorpdfstring{$s<\mathrm{lpf}(k)$}{s < lpf(k)}}

Theorem \ref{thm:B} used $k$ prime. The sharp general criterion appears to be:

\begin{proposition}\label{prop:lpf}
$E^0_{s,s}(k)=\mathrm{P}(s,k)$, i.e.\ characteristic zero contributes no
non-trivial solution, if and only if $s<\mathrm{lpf}(k)$. This is proved for $k$
prime (Lemma \ref{lem:prime-k}) and for $k=q^e$; for general $k$ the ``only if''
is elementary and the ``if'' is verified in $144$ cells with no exception.
\end{proposition}
\begin{proof}[Proof for $k=q^e$]
A multiple $f=\Phi_{q^e}g$ of degree $<q^e$ has $\deg g<q^{e-1}$, so the $q$
blocks of $\Phi_{q^e}(x)=\sum_{j<q}x^{jq^{e-1}}$ do not overlap and the
coefficient vector of $f$ is that of $g$ repeated $q$ times. Hence the positive
mass of $f$ is $q$ times that of $g$, so at least $q$; and $f(1)=q\,g(1)=0$
forces $g$ to carry both signs, so it is nonzero on both sides.
\end{proof}

For general $k=q\,k'$ with $q=\mathrm{lpf}(k)$ the tensor decomposition
$\ZZ[\zeta_k]=\ZZ[\zeta_q]\otimes\ZZ[\zeta_{k'}]$ reduces the relation to
$\sum_{\alpha\in\ZZ/q}c_\alpha\zeta_q^\alpha=0$ with
$c_\alpha\in\ZZ[\zeta_{k'}]$, whence all $c_\alpha$ are equal to a common $c$.
If $c=0$ the statement follows by induction on the number of prime factors, since
every $\mathrm{lpf}(k')>q>s$. If $c\neq0$ then no $\alpha$-class is empty on both
sides; if every class contains a term of the first tuple then $s\ge q$, and
otherwise Lam--Leung applied to the resulting vanishing sum over $\mu_{k'}$ gives
a class of size $\ge\mathrm{lpf}(k')>q$. The remaining gap is a factor $2$ in the
last branch, which is why we record the general case as verified rather than
proved. Cells checked: $54$ (mixed $k\le41$), $29$ (prime powers
$k\le81$), $61$ (composite $k\le77$ with two prime factors) --- $144$ in total,
zero exceptions.

\subsection{Lemma F: the exact quadratic phase, and what the barrier really is}
\label{sub:lemmaF}

T7 asserts only $|\Tr(T(v)U^J)|=1$ for $J\not\equiv0\ (K)$. That modulus-only
statement is the source of the exponential error $C^r/m$ in the moment formula T8
and hence of the T18 barrier. It can be sharpened to an identity.

Write $\sigma(m,v)=m_2v_1-m_1v_2$ for the symplectic form, so that
$T(m)T(v)=\omega^{2^{-1}\sigma(m,v)}T(m+v)$ and
$T(m)T(v)T(m)^{-1}=\omega^{\sigma(m,v)}T(v)$, and put $B_J=(A^J)^{\mathsf T}$,
which is invertible with $B_J-\mathbf1$ invertible for $J\not\equiv0$.

\begin{lemma}\label{lem:F}
For $J\not\equiv0\pmod K$ set $m(v)=-(B_J^{-1}-\mathbf1)^{-1}v$, a linear
function of $v$. Then
\[
  \Tr\bigl(T(v)U^J\bigr)=\Tr\bigl(U^J\bigr)\,\omega^{\,Q_J(v)},
  \qquad
  Q_J(v)=\sigma\bigl(m(v),v\bigr)-\tfrac12\,\sigma\bigl(B_J^{-1}m(v),m(v)\bigr),
\]
and $Q_J$ is a \emph{homogeneous} quadratic form in $v$ (no linear or constant
term), with $|\Tr(U^J)|=1$.
\end{lemma}

\begin{proof}
Conjugating by $T(m)$ and using
$T(m)U^JT(m)^{-1}=\omega^{-\frac12\sigma(B_J^{-1}m,m)}
 U^JT\bigl((B_J^{-1}-\mathbf1)m\bigr)$ gives
\[
  \Tr(T(v)U^J)=\omega^{\sigma(m,v)-\frac12\sigma(B_J^{-1}m,m)}
   \Tr\bigl(T(v)T((B_J^{-1}-\mathbf1)m)U^J\bigr),
\]
and choosing $m$ with $(B_J^{-1}-\mathbf1)m=-v$ collapses the right-hand trace to
$\Tr(U^J)$. Homogeneity is clear since $m(v)$ is linear.
\end{proof}

Verification: the quadratic-form structure was \emph{determined} from six samples
by linear algebra over $\ZZ/N$ and then tested out of sample --- $5000$ checks
over five $(A,N)$ and $J\le5$, \textbf{zero mismatches}, and every fitted linear
and constant coefficient is exactly $0$. The closed form above then matches the
fitted coefficients under one universal sign convention at all $16$ triples
$(A,N,J)$ tested.

\begin{corollary}\label{cor:barrier}
The error term of the moment formula is
\[
  \theta_r=\frac1{mK}\sum_{J\not\equiv0}\lambda^{-J}\Tr(U^J)
   \sum_{v}\widehat{(\bar V^{\,r})}(v)\;\omega^{\,Q_J(v)} ,
\]
where $\widehat{(\bar V^{\,r})}$ is the Weyl coefficient vector of $\bar V^{\,r}$.
Bounding the inner sum by its $\ell^1$ norm reproduces T8's $C^r/m$; any
square-root cancellation in this explicit \emph{theta-like} sum against a
homogeneous quadratic phase improves the T18 threshold directly.
\end{corollary}

This replaces ``a local law is needed'' by a concrete analytic target: cancellation
in $\sum_v a_v\,\omega^{Q_J(v)}$ for the specific coefficient vectors
$a=\widehat{(\bar V^{\,r})}$, $r$ up to $\asymp\log m$. It is the sharpest
formulation of the obstruction that we can give, and it is the natural place for
the next attack: the phases are Gauss-type, the coefficient vector is explicit
and sparse for small $r$, and the only reason the $\ell^1$ bound is used at
present is that no structural information about $a$ has been exploited.

\subsection{Theorem G: the moment route is provably closed}\label{sub:thmG}

Lemma \ref{lem:F} localises the obstruction in an explicit theta-like sum; the
next question is whether that sum enjoys cancellation. It does not, and the
reason is structural.

\begin{theorem}\label{thm:G}
Let $V$ have Weyl coefficients $c$ with $\|c\|_1=O(1)$ and let
$\bar V=K^{-1}\sum_{j<K}U^jVU^{-j}$ be the $K$-step Birkhoff average, whose Weyl
coefficient vector is $\bar c$. Then
\[
  \|\bar c\|_1\asymp\|c\|_1,\qquad
  \|\bar c\|_2\asymp K^{-1/2}\|c\|_2,\qquad\text{so}\qquad
  \frac{\|\bar c\|_1}{\|\bar c\|_2}\asymp\sqrt K\,\frac{\|c\|_1}{\|c\|_2},
\]
because the Birkhoff average spreads each frequency $n$ over its full
$A^{\mathsf T}$-orbit with weight $K^{-1}$, and distinct orbits do not interfere.
Consequently, with $a_r$ the $v=0$ Weyl coefficient of $\bar V^{\,r}$ (i.e.\ the
classical moment) and $\theta_r$ the error of T8,
\[
  \frac{|\theta_r|}{|a_r|}\ \lesssim\ \frac{1}{m}
   \Bigl(\frac{\|\bar c\|_1}{\|\bar c\|_2}\Bigr)^{r}
   \asymp\frac{(c\sqrt K)^r}{m},
\]
which is $o(1)$ only for $r\lesssim 2\log m/\log K$. In particular no observable
with $O(1)$ Fourier support can reach $r\asymp\log m$: the dilution factor
$K^{-1}$ is forced, and an $A$-invariant observable supported on a single orbit has
$\|\bar c\|_1/\|\bar c\|_2\asymp\sqrt K$ with no freedom left.
\end{theorem}


Whether the crude $\ell^1$ bound in Theorem \ref{thm:G} could be beaten is a
question about the theta-sum of Corollary \ref{cor:barrier}, and it is decided
empirically. For the single-orbit observable we computed the Weyl coefficient
vector $a=\widehat{(\bar V^{\,r})}$ exactly by the $O(N^2\log N)$ Weyl transform
and evaluated every $S_J(r)=\sum_va_v\omega^{Q_J(v)}$:

\begin{center}\footnotesize
\begin{tabular}{@{}lrrrr@{}}
\toprule
& $r=4$ & $r=5$ & $r=6$ & $r=8$\\
\midrule
$\max_J|S_J|/\|a\|_1$, $N=41$   & $0.448$ & $0.525$ & $0.350$ & $0.324$\\
$\max_J|S_J|/\|a\|_1$, $N=101$  & $0.287$ & $0.331$ & $0.239$ & $0.221$\\
$\max_J|S_J|/\|a\|_1$, $N=71$   & $0.310$ & $0.312$ & $0.276$ & $0.245$\\
$\max_J|S_J|/\|a\|_2$ (all)     & $7.9$--$14.5$ & $9.9$--$18.2$ & $7.5$--$14.6$ & $7.6$--$14.9$\\
\bottomrule
\end{tabular}
\end{center}

\noindent
So $|S_J|$ sits a \emph{constant factor} $\approx3$ below the $\ell^1$ bound and
one to two orders of magnitude \emph{above} $\|a\|_2$: there is no square-root
cancellation to harvest. The $C^r$ growth in T8 is therefore genuine and not an
artefact of a lossy estimate, and:

\begin{corollary}
The moment method cannot reach the mean-spacing scale for this family, for any
fixed observable of bounded Fourier complexity. Any proof of BGS here must be
resolvent-based, not moment-based.
\end{corollary}

This is the sharpest negative statement in the programme and it is a useful one:
it closes an entire route, explains why Theorem \ref{thm:A} succeeds (the exact
covariance identity bypasses moments altogether rather than pushing them), and
leaves a single well-posed question --- a self-consistent equation for
$(V_\lambda-z)^{-1}$ at $\operatorname{Im}z\asymp m^{-\delta}$ --- as the whole
remaining gap.

\subsection{A three-way subspace control}\label{sub:control}

An attempt to push the desymmetrised measurement to $m=1674$ ($p=95419$, $K=57$)
was set up with target $m=1280$, i.e.\ a \emph{proper} subspace of the
$1674$-dimensional eigenspace; a proper subspace of $\HH_\lambda$ is not parity
invariant, so what it measured is not a desymmetrisation. The diagnostic that
exposed this was already in the output: the off-parity block of the compression
was $1.4$--$1.6\cdot10^{-2}$, against $\le2.2\cdot10^{-15}$ in the genuine runs.
Recorded as a control, it is a decisive one:

\begin{center}\footnotesize
\begin{tabular}{@{}lcc@{}}
\toprule
subspace & $\langle\tilde r\rangle$ & reading\\
\midrule
Haar-random subspace of $\CC^N$ & $0.593$--$0.598$ & $\GUE$, $z_{\GUE}=-0.5\dots-1.1$\\
random subspace inside $\HH_\lambda$ ($1280$ of $1674$) & $0.5963\pm0.0034$
 & $\GUE$, $z_{\GUE}=-0.45$, $z_{\GOE}=+9.09$\\
\textbf{full} $\HH_\lambda$, parity blocks & $0.529$--$0.542$
 & $\boldsymbol{\GOE}$, $z_{\GOE}=-0.9\dots+0.4$\\
\bottomrule
\end{tabular}
\end{center}

\noindent
The orthogonal class therefore lives on the \emph{full} eigenspace, where the
antiunitary $\Theta_u$ of Proposition \ref{prop:Esquare} acts; restricting to a
non-invariant subspace returns $\GUE$. This is an independent confirmation of the
mechanism of Theorem \ref{thm:E}, obtained by accident. The largest bona fide
desymmetrised measurement is now $m=1674$ (at $p=95419$, $K=57$: the parity
blocks give $\langle\tilde r\rangle=0.5290\pm0.0022$ and
$0.5344\pm0.0042$ against the $\GOE$ reference $0.5314\pm0.0125$, i.e.
$z_{\GOE}=-0.19,+0.24$ at $837$ levels per block, $9\sigma$ from $\GUE$);
the prior limitation to $m=1280$ was an 8~GB memory ceiling, re-probed at
16~GB.


\section{The frequency ensemble: the route Theorem G does not close}
\label{sec:ensemble}

Theorem \ref{thm:G} closes the moment route for a \emph{fixed} observable. But BGS
in this family is naturally a statement about the ensemble
$V_n=T(n)+T(-n)$, $n\in(\ZZ/N)^2$, and averaging over $n$ changes the estimate
qualitatively.

By the Birkhoff reduction $P_\lambda T(v)P_\lambda=W_vP_\lambda$ with
$W_v=K^{-1}\sum_jT(A^{\mathsf Tj}v)$, the Weyl relations collapse every product:
\[
  \tfrac1m\Tr(M_n^r)=K^{-r}\!\!\sum_{\epsilon,j}\!\varphi(\epsilon,j)
   \mathbf 1\{L(\epsilon,j)n=0\}
  +\frac1{mK^r}\sum_{\epsilon,j}\varphi\,E\bigl(L(\epsilon,j)n\bigr),
  \quad L(\epsilon,j)=\sum_i\epsilon_i(A^{j_i})^{\mathsf T}.
\]

\begin{proposition}[the rank-zero condition is two Chebyshev congruences]
Write $A^j=U_jA-U_{j-1}\mathbf1$ with $U_0=0$, $U_1=1$, $U_{j+1}=tU_j-U_{j-1}$,
$t=\tr A$, $U_{-1}=-1$. Then
\[
  L(\epsilon,j)\equiv0\ (N)\iff
  \sum_i\epsilon_iU_{j_i}\equiv0\ \text{and}\ \sum_i\epsilon_iU_{j_i-1}\equiv0
  \pmod N,
\]
two signed vanishing sums of a lacunary integer sequence --- the matrix analogue of
the characteristic-zero rigidity of Theorem \ref{thm:B}. Certified on $1600$
random $(\epsilon,j)$ over four $(A,N)$ including $K=200$, zero mismatches.
\end{proposition}

\begin{proposition}[averaging over $n$ gains a factor $N$]\label{prop:gain}
By Lemma \ref{lem:F} the error carries the phase $\omega^{Q_J(Ln)}$, a quadratic
form in $n$, so $\mathbb E_n[\omega^{Q_J(Ln)}]=N^{-2}\cdot(\text{binary Gauss
sum})=O(1/N)$ whenever $Q_J\circ L$ is nondegenerate. Hence the error term, of
size $2^rK/m$ for fixed $n$, becomes $O(2^r/(mN))$ after averaging, and with
$K=N^\theta$ the reachable moment order moves from $r<2/\theta$ to $r<4/\theta$.
\end{proposition}

The cancellation that Theorem \ref{thm:G} measured to be \emph{absent} in $v$ is
therefore \emph{present} in $n$. What must then be controlled is the count
$D_r(N)=\#\{(\epsilon,j):L=0\ (N)\}$, and with $r=2s$ the number of balanced words
over $\{(j,\pm)\}$ is $r![x^r]I_0(2x)^K=\binom{2s}{s}\mathrm{P}(s,K)$ with
$\mathrm{P}$ the pairing count of Lemma \ref{lem:pair}. Measured:

\begin{center}\footnotesize
\begin{tabular}{@{}rrrrrrr@{}}
\toprule
$N$ & $K$ & $\varepsilon$ & $r$ & $D_r(N)$ & balanced-word count & excess\\
\midrule
$3571$ & $17$ & $+1$ & $2$ & $34$ & $34$ & $\mathbf 0$\\
$3571$ & $17$ & $+1$ & $4$ & $3366$ & $3366$ & $\mathbf 0$\\
$3571$ & $17$ & $+1$ & $6$ & $545020$ & $538900$ & $+1.1\%$\\
$3571$ & $17$ & $+1$ & $8$ & $122637830$ & $117173350$ & $+4.5\%$\\
$9349$ & $19$ & $+1$ & $4$ & $4218$ & $4218$ & $\mathbf 0$\\
$9349$ & $19$ & $+1$ & $8$ & $193271610$ & $186398170$ & $+3.6\%$\\
\bottomrule
\end{tabular}
\end{center}

\noindent
The count is \emph{exact} at $r=2,4$, and the excess at $r=6,8$ is already present
over $\ZZ$ ($5040$ of $6120$; $4.45\cdot10^6$ of $5.46\cdot10^6$): it is
combinatorial, produced by the Chebyshev recurrence
$U_{j+1}+U_{j-1}=tU_j$ written with repeated terms, and not a mod-$N$ accident.
For $\varepsilon=-1$ ($A^k\equiv-\mathbf 1$, $K=2k$) the antiperiodicity $U_{j+k}=-U_j$ makes the vectors $(1$ at $a$, $1$ at $a+k)$ --- the mod-$N$ relations $\mathbf 1 + x^k$ --- a MODULAR family already at $r=2$: the excess is dominated by explicitly parametrised antiperiodic classes (($x^a\pm x^{a+k}$)-products), the hierarchy opposite to the $\varepsilon=+1$ case, and the naive factor $2^{r/2}$ of the first census is retired (run 78: the corrected decomposition reproduces the independent meet-in-the-middle counts $D_2=136$, $D_4=53856$, $D_6=34881280$ at $K=34$ exactly, and $D_4=2229120$ at $K=216$ via the explicit family formulas).

\begin{remark}
The whole remaining gap is now: bound the Chebyshev $\pm1$-relation count $X_r$ for
$r\asymp\log m$, then run a resolvent argument on the ensemble moments. Both halves
are explicit. Note also that a signed sum of $r$ terms of a sequence bounded by
$\varphi^K\asymp N$ lies in $[-rN,rN]$, so the mod-$N$ condition reads
$\sum\epsilon_iU_{j_i}=cN$ with only $O(r)$ admissible values of $c$ --- the
arithmetic part of the problem is finite and small.
\end{remark}

% The cyclic shape decomposition of D_r
\subsection{The cyclic shape decomposition of \texorpdfstring{$D_r$}{Dr}}\label{sub:thmH}

The remark above asked to bound $X_r$.  This subsection solves its
\emph{structural} half exactly.  Group the words $(\epsilon_i,j_i)_{i\le r}$,
$j\in[0,K)$, by their \emph{collected coefficient vector}
$c_j=\sum_{i:\,j_i=j}\epsilon_i$.  The signature depends only on $c$, and over
the integers the rank-zero condition becomes a divisibility statement:

\begin{lemma}[collected vectors]\label{lem:coll}
$\sum_i\epsilon_iA^{j_i}=0$ over $\ZZ$ $\iff$ $m_t\mid p_c$ in $\ZZ[x]$, where
$p_c(x)=\sum_jc_jx^j$, because $A^j=U_jA-U_{j-1}\mathbf 1$ and $\{A,\mathbf 1\}$
spans $\mathbb Q[A]$.  The number of length-$r$ words with collected vector $c$ is
\[
  \operatorname{wc}(c,r)\;=\;r!\,[y^r]\prod_{j<K}I_{|c_j|}(2y),
  \qquad [y^{2k}]\,I_a(2y)=\tfrac1{k!(k+a)!},
\]
where cancelled pairs (the case $\sum_j|c_j|<r$) are carried by the $I_0$
factors, and parity forces $\sum_j|c_j|\equiv r \pmod 2$.
\end{lemma}

Call a \emph{shape} a support-normalised nonzero multiple of $m_t$, and let
$\mathcal S(R)$ be the shapes of $\ell^1$ norm exactly $R$.  For
$\varepsilon=+1$ designs ($A^K\equiv\mathbf 1\ (N)$):

\begin{theorem}[cyclic shape decomposition]\label{thm:H}
\[
  D_r(N)\;=\;\binom{2s}{s}\mathrm P(s,K)\;+\;
  2\!\!\sum_{\substack{R\equiv r\ (2)\ R\le r}}\;\sum_{\sigma\in\mathcal S(R)}
  \;\sum_{a\in\mathbb Z/K}\operatorname{wc}\bigl(\mathrm{dense}(a.\sigma\bmod K),K,r\bigr)
  \;+\;M_r(N),
\]
where $a.\sigma$ places $\sigma$ at positions $a,\dots,a+\mathrm{width}-1$
reduced into the cyclic window $[0,K)$, and $M_r(N)$ collects the genuinely
modular vectors: signature $0\ (N)$ with no cyclic rotation a $\ZZ$-multiple.
\end{theorem}

\begin{proof}
Grouping by $c$ is exact.  The atom signatures are $K$-periodic mod $N$, and the
Chebyshev shift identity $U_{i+a}=U_aU_{i+1}-U_{a-1}U_i$, applied to the two
vanishing sums, shows that every cyclic shift of a vanishing configuration again
vanishes mod $N$; injectivity of $(\sigma,a,\pm)$ on support-normalised shapes is
immediate, and the remainder is $M_r$ by definition.
\end{proof}

The formula is certified \emph{to the last digit} against the table above, and
against a fresh meet-in-the-middle at $r=5$: all $12$ comparison rows ($24$
numbers) pass with zero mismatches, and the full
census of all collected vectors with $0<\|c\|_1\le6$ reproduces $D_6$
vector-by-vector.  Three consequences deserve emphasis.

\emph{(i) There is only one excess mechanism.}  The $r=6$ census shows that the
\emph{entire} mod-$N$ accident class ($1080$ words at both $K=17$ and $K=19$)
consists of the cyclic \emph{wraps} of the shape $(1,-2,-2,1)$ across the period
boundary; $M_r(N)=0$ for $r\le8$.  The R74 reading --- combinatorial excess over
$\ZZ$ plus modular accidents --- is superseded: both are the same shapes, shifted
and wrapped.

\emph{(ii) Fixed-order convergence.}  The excess is $O(K\cdot\text{shape mass})$
while the main term is $\asymp K^s$, so $X_r/\text{main}=O_r(K^{1-s})\to0$ at
every fixed $r\ge4$.  Measured exactly at $r=6$: relative excess $1.123\%$
($K=17$) vs $0.892\%$ ($K=19$), ratio $1.259$ against the $K^{-2}$ prediction
$1.249$.

\emph{(iii) Odd orders are pure arithmetic.}  $D_5(\ZZ)=2(K-2)\cdot20$ and
$D_5(N)=2K\cdot20$ exactly ($600/680$ at $K=17$; $680/760$ at $K=19$): the
ensemble fifth moment is $\asymp K^{-4}$, nonzero, and produced by the single
shape $(1,-3,1)=m_t$.  The vanishing odd moments of $\GUE$ are recovered
asymptotically with an explicitly arithmetic deviation.

The next section attaches the phases.  With the product law
$T(m)T(v)=\omega^{\sigma(m,v)/2}T(m+v)$ the fixed-$n$ moment carries the
closed-form phase $\Phi=\frac12\sigma(n,C(\epsilon,j)n)$ with
$C=\sum_{i<k}\epsilon_i\epsilon_k A^{(j_k-j_i)\mathsf T}$, and the formula is
certified \emph{end-to-end} against direct matrix computation at
$(N,K,m)=(199,11,18)$ for $r=2,\dots,5$ and six values of $n$: $24/24$
agreements, max error $8\cdot10^{-14}$.  Three level facts follow: the
pairing words carry $\varphi\equiv1$ identically; $E(0)=0$; and the
$n$-average of a rank-zero word is a two-dimensional quadratic Gauss sum
$\mathcal G(C)=N^{-2}\sum_n\omega^{\sigma(n,Cn)/2}$ with
$|\mathcal G|\le N^{-1/2}$, saturation $|\mathcal G|=N^{-1}$ for
nondegenerate $C$.  Measured at $(3571,17)$: the $680$ $m$-shape words carry
four distinct $C$ matrices, all nondegenerate, $|\mathcal G|=N^{-1}$ exactly,
and $G_5=\sum\mathcal G(C)=0.190423-0.000000i$ --- real to machine precision.
The rank-zero part of the ensemble fifth moment is therefore
$K^{-5}G_5=+1.34\cdot10^{-7}$: the ensemble excess is the Theorem
\ref{thm:H} counting excess suppressed by an extra $N^{-1/2}$ (measured
$N^{-1}$).  A vectorised census at $r=10$ (cross-validated exactly at
$r=6,8$) gives $3\,860\,398\,080$ excess words $=3\,107\,822\,760$
$\ZZ$-multiples $+\,752\,575\,320$ wraps $+\,\mathbf 0$ genuinely modular:
through $r=10$ the cyclic shape sum remains the whole excess, and the
relative excess at $K=17$ grows through $1.14\%$, $4.66\%$, $12.15\%$ at
$r=6,8,10$ --- the empirical uniformity ceiling that any bound on the shape
mass must reproduce.

The shape census for $m_3$ (width $\le16$) is $S(R)=1,1,1,3,8,35,64,105$ for
$R=5,\dots,12$ (a width-$14$ pass gives $31,56,93$ at $R=10,11,12$: the wide
sparse shapes are real, so the counts grow with the allowed width); the first shape being $(1,-3,1)=m_t$ itself ($\ell^1$ norm
$t+2$); for $m_6$ nothing at all exists below $R=8$ besides $(1,-6,1)$.  The
remaining gap of Proposition \ref{prop:gain} is therefore reduced to the
$\varphi$-weights per shape, a uniform-in-$r$ bound on the shape mass (the
Littlewood-multiple counting problem for $x^2-tx+1$), and the resolvent step.

\subsection{The generating function evaluated; a correction to
Theorem \texorpdfstring{\ref{thm:H}}{H}}\label{sub:r80}

This section evaluates the generating function of \S\ref{sub:thmH} \emph{exactly} and, in doing
so, exposes a correction to Theorem \ref{thm:H} that the $r\le8$
certification window was structurally blind to.

\emph{The $J$-weighted automaton.}  Each accepting walk of the
reversed-normalized automaton (the support-normalised multiples of
$m_t$, read top-down) carries the letter weight
$J_c(y)=I_c(2y)/I_0(2y)$, internal zero gaps carrying $J_0=1$.  With
\[
  \Psi(y)=\sum_{\text{walks}}\prod_{\text{letters}}J_{|c|}(2y),
  \qquad
  \Psi_w(y)=\sum_{\text{walks}} w\prod_{\text{letters}}J_{|c|}(2y)
\]
($w$ the walk width), the identity
$I_0(2y)^K\prod_{\text{letters}}J=\prod_{j<K}I_{|c_j|}(2y)$ (the dense
$K$-window product, cancelled pairs included) turns Theorem \ref{thm:H}'s
counts into exact coefficient extractions valid for every $r\le r_{\max}$
and $K\le K_{\max}$ simultaneously:
\[
  X^{\ZZ}_r(K)=2r!\,[y^r]\,I_0(2y)^K\bigl((K{+}1)\Psi-\Psi_w\bigr),
  \qquad
  \operatorname{triv}(r,K)=r!\,[y^{2s}]\,I_0(2y)^K\ (r=2s),
\]
the rotation-inclusive count being $2Kr!\,[y^r]I_0(2y)^K\Psi$.  Note the
exact rank-zero count $\operatorname{triv}(r,K)=r![y^{2s}]I_0(2y)^K$ (it
gives $\operatorname{triv}(6,17)=538\,900$): the display
$\binom{2s}{s}\mathrm P(s,K)$ in Theorem \ref{thm:H} is shorthand only.
Certified against the census at every available row: $r\le8$ exact
($\ZZ$/full/wrap at $K=17,19$), $X^{\ZZ}$ exact through $r=12$ including
the padded shapes, $\mathrm S_{\mathrm{win}}(R;17)$ $8/8$ against the
reversed-normalized automaton, and an independent float implementation agreeing to
$2.2\cdot10^{-15}$ (\texttt{r80\_jweight\_dp.py}, \texttt{r80\_julia\_bigdp.jl}).

\emph{Wrap coincidences.}  The rotation-inclusive count equals the census
exactly through $r=8$ and \emph{exceeds} it from $r=10$.  The
$(\sigma,a,\pm)$ parametrisation of Theorem \ref{thm:H} is not injective
once padded shapes exist: a wrapping rotation of a shape can equal another
shape as a signed $K$-vector.  The mechanism is automatic --- if
$\sigma_A=\pm x^a\sigma_B \pmod{x^K-1}$ with both $\sigma_{A,B}$
$m_t$-multiples, then the difference is $(x^K-1)h$ with $m_t\mid h$, since
$\gcd(m_t,x^K-1)=1$ over $\mathbb Q$ (the roots of $m_t$ are real units,
not roots of unity) and Gauss's lemma lifts the divisibility to
$\mathbb Z[x]$.  The first collisions occur at $\ell^1$ norm $10$, in the
padded family $m_t(1+x^d)$ (e.g.\ at $K=17$,
$\mathrm{rot}_{14}\,m_t(1+x^3)=m_t(1+x^{14})$ as vectors).  The exact
excess of the rotation-inclusive count over the census,
$\sum_v(\mathrm{mult}(v)-1)\operatorname{wc}(v,r)$ over signed vectors $v$,
is $41\,126\,400$ ($r=10$), $1\,244\,073\,600$ ($r=11$) and
$97\,037\,740\,800$ ($r=12$) at $K=17$ --- both checkable census deltas
reproduced to the last digit (\texttt{r80\_wrap\_collisions.py}).  The
correct statement of Theorem \ref{thm:H} is therefore at the level of
\emph{distinct} collected vectors,
$D_r-\operatorname{triv}-M_r=\sum_{\text{distinct }c}\operatorname{wc}(c,r)$,
with the shapes of width $\le K$; the census reading, $M_r=0$ through
$r=12$, is unaffected.

\emph{The shape growth.}  The count DP (\texttt{r80\_julia\_bigdp.jl},
exact integer arithmetic) gives, for the first time without a width cap:
dense shapes (width $\le$ norm) $S(30)=301\,160\,711$, $S(50)=7.6\cdot
10^{15}$, $S(80)=1.17\cdot10^{18}$, with $S(R)^{1/R}$ still rising ---
$2.078$ at $R=50$, $2.179$ at $R=80$; the width-capped ($\le80$) curve
peaks at $2.2499$ ($R=47$) and turns down, the $C^R\operatorname{poly}(R)$
signature.  The growth constant of the language therefore satisfies
$2.18<C_{\mathrm{dense}}\le3$, the upper bound being the elementary
signed-compositions count.  Pinning $C_{\mathrm{dense}}$ (Perron analysis
of the co-reachable automaton core, or a Frougny--Steiner transfer) is the
named remaining arithmetic problem of the uniformity step.


\subsection{The observable family is a one-parameter arithmetic
curve}\label{sub:r82}

The family of compressed observables studied throughout this programme,
\[
  V_n \;=\; P_+\bigl(T(n)+T(-n)\bigr)P_+ ,
  \qquad n\in(\ZZ/N)^2\setminus\{0\},
\]
carries an exact rigidity that reduces it from $(N^2-1)/K$ nominal members
(the $A^{\mathsf T}$-orbits) to at most $N$.  Let $T\le SL_2(\ZZ/N)$ be the
centraliser of $A$ -- the group whose Weil operators are the \emph{Hecke
operators} of Kurlberg and Rudnick \cite{KR2000} -- and let
\[
  Q(n)\;=\;n_1^2-n_1n_2-n_2^2 \pmod N
\]
be its invariant quadratic form (discriminant $t^2-4=5$ at $t=3$).

\begin{theorem}[torus rigidity]\label{thm:torus}
Let $N$ be an odd prime with $\gcd(N,t^2-4)=1$.  If $Q(n)=Q(n')$ then
$V_{n'}=U_gV_nU_g^{*}$ for an explicit $g=aI+bA\in T$; in particular
$\operatorname{spec}(V_n)=\operatorname{spec}(V_{n'})$.  Consequently the
family has at most $N$ distinct spectra, one for each value of $Q$; each
nonzero level set of $Q$ is a single $T$-orbit containing exactly $|T|/K=m$
distinct $A^{\mathsf T}$-orbits, and
\[
  (N\mp1)\,m \;+\; \frac{|\{Q=0\}\setminus\{0\}|}{K}\;=\;\frac{N^2-1}{K}.
\]
\end{theorem}

\begin{proof}
Identify $(\ZZ/N)^2$ with $R_A=(\ZZ/N)[x]/(x^2-tx+1)$ as a free rank-one
$R_A$-module.  Then $Q$ is a scalar multiple of the norm form of $R_A$ and $T$
is its group of norm-one units, so $T$ acts simply transitively on each level
set $\{Q=q\}$, $q\neq0$; in particular $n'=gn$ for some $g\in T$.  Since $g$
commutes with $A$, the Weil operator $U_g$ commutes with $U_A$ up to a scalar
and therefore preserves every $U_A$-eigenspace; since $-I$ is central, $U_g$
commutes with the parity operator and preserves each parity block.  The Weil
intertwining property $U_gT(n)U_g^{*}=T(gn)$ then gives
$U_gV_nU_g^{*}=P_+(T(gn)+T(-gn))P_+=V_{n'}$.
\end{proof}

Theorem \ref{thm:torus} is elementary given \cite{KR2000}, but it settles a
question this programme had left open, and it constrains the ensemble route
sharply.

\begin{corollary}\label{cor:curve}
The $n$-ensemble $\{V_n\}$ is a deterministic one-parameter curve
$q\mapsto V_q$, each member repeated $m$ times, and every average over
$n\in(\ZZ/N)^2$ equals the corresponding average over $q\in\FF_N^{*}$.
Hence \textup{(i)} the ensemble has one intrinsic fluctuation direction, so a
single dominant coherent mode and positively correlated bulk window counts are
automatic; \textup{(ii)} no i.i.d.\ $\GOE$ statement can hold for $\{V_n\}$;
\textup{(iii)} averaging over $n$ buys concentration in one arithmetic scalar
only and therefore cannot substitute for a per-system local law.
\end{corollary}

Part \textup{(iii)} closes, for this family, the last route in which an
ensemble average might have replaced a local law; together with the
moment-route obstruction it leaves the single-matrix problem: for fixed $q$,
prove $\GOE$ local statistics for $V_q$.  In the Hecke basis $\{f_j\}_{j\in
\ZZ/m}$ (the $T$-character basis) equivariance forces
$\langle f_j|T(n)|f_k\rangle=c_{jk}(q)\,\zeta^{(j-k)s(n)}$ with
$\zeta=e^{2\pi i/m}$ and $s(n)$ the discrete logarithm of the torus element
carrying a base point to $n$; so
$V_n=D_{s(n)}(C_q+C_q^{*})D_{s(n)}^{*}$ with $D_s$ diagonal, and the entire
arithmetic content is one matrix $C_q=(c_{jk}(q))$ of explicit exponential
sums per level set.

\emph{The modulation, measured exhaustively.}  With $q$ as the parameter the
curve is computable over \emph{all} level sets.  At $(N,K)=(3571,17)$,
$m_+=105$, over all $3570$ values of $q$: $\operatorname{Tr}(V_q^2)/m_+$ has
mean $0.117600$ against the arithmetic value $N/(K^2m_+)=0.117658$ and
relative spread $1.22\%$, while
$\mu_4(q)\in[1.2246,2.8573]$ with mean $1.9694$ and standard deviation
$0.2350$ (the semicircle value is $2$; the arcsine value $3/2$ is excluded).
The corresponding matched $\GOE$ reference has $\operatorname{std}(\mu_4)
=0.027$, so the macroscopic shape of $V_q$ does not self-average.  Writing
$A^j=u_jI+v_jA$, so that $v_j=1,3,8,21,55,144,377,987,\dots$ (the even-index
Fibonacci numbers at $t=3$), the additive Fourier transform in $q$ of the raw
central traces is supported on those coefficients: the $16$ frequencies
$\pm v_j$, $1\le j\le 8$, carry $90\%$ of the variance of
$\operatorname{Tr}(V_q^4)$ and $\operatorname{Tr}(V_q^6)$, and a
least-squares fit on exactly those $(K-1)/2$ frequency pairs gives
$R^2=0.9426$, $0.9276$ and, for $\mu_4(q)$, $R^2=0.98655$.  Enlarging the
frequency set to $71$ frequencies moves $R^2$ only to $0.9472$/$0.9872$: the
comb is saturated.  The shape is that of a weighted sum of $(K-1)/2$
oscillations, matching the form conjectured by Kurlberg and Rudnick for the
matrix elements of the desymmetrised cat map \cite{KR2003}; the coefficients
are fitted here, not derived.

\emph{Where the excess lives.}  In the basis diagonalising linear eigenvalue
statistics, $L_k=\sum_iT_k(\lambda_i/\sqrt2)$, the variance excess over the
matched $\GOE$ reference is confined to \emph{even} modes: the ratios at
$k=4,6,8,10$ are $81,71,78,100$ at $m_+=105$ and $375,306,302,352$ at
$m_+=246$, against $1.0,1.6,2.8,5.8$ and $0.5,2.9,5.2,10.8$ at
$k=3,5,7,9$; a Haar control on the same operators gives $0.3$--$0.6$
throughout.  Removing the component of the bulk count vectors explained by
$k\le8$ drops the count-variance ratio from $1.89$ to $1.12$ ($m_+=105$) and
from $3.82$ to $1.46$ ($m_+=246$).  The excess is not an edge effect: the
largest normalised eigenvalue sits at $1.35$ bulk edges on average against
$\GOE$'s $1.39$, and trimming the eight most deviant eigenvalues leaves the
$\mu_4$ variance ratio at $52$ (from $69$).  This is exactly the signature of
a one-parameter family of symmetric macroscopic densities.


\subsection{The family in closed form, and BGS at large
\texorpdfstring{$m$}{m}}\label{sub:r82b}

Theorem \ref{thm:torus} says the family is a curve; the following says what the
curve is.

\begin{theorem}[twisted Jacobi sums]\label{thm:jacobi}
Let $N$ be an odd prime, split for $A$ (so $t^2-4$ is a square mod $N$).
Conjugating inside $SL_2(\ZZ/N)$, take $A'=\operatorname{diag}(a,a^{-1})$ with
$a^2-ta+1=0$ and $\operatorname{ord}(a)=K$.  Then the Hecke basis is the set of
multiplicative characters, $f_\chi(x)=\chi(x)/\sqrt{N-1}$ for $x\neq0$ and
$f_\chi(0)=0$, with $U_uf_\chi=\left(\tfrac uN\right)\chi(u)^{-1}f_\chi$; a
$U_{A'}$-eigenspace is $\{\chi:\chi(a)=\text{const}\}$ of dimension $(N-1)/K$;
parity splits it into even and odd characters, $m_\pm=(N-1)/(2K)$; and, for
$n_1\neq0$, with $q=n_1n_2$,
\[
  \langle f_\chi|T(n)|f_\psi\rangle
  =\ep(-q/2)\,\overline{\chi}(n_1)\psi(n_1)\,
   \frac{J(\chi,\psi;q)}{N-1},
  \qquad
  J(\chi,\psi;q)=\sum_{x\neq0,1}\overline{\chi}(x)\,\psi(x-1)\,\ep(qx).
\]
\end{theorem}

The hypothesis $n_1\neq0$ is load-bearing: on the line $n_1=0$ the printed
right-hand side vanishes identically ($\overline{\chi}(0)=0$), while the true
matrix element $\langle f_\chi|T((0,n_2))|f_\psi\rangle$ is a Gauss-sum
diagonal object of size $O(N^{-1/2})$, nonzero in general (measured up to
$0.218$ at $N=23$, decaying as $N^{-1/2}$;
\texttt{review/referee/ref03\_theoremII\_jacobi\_derivation.py}).  Every
certified consequence below uses level sets with $q=n_1n_2\neq0$, where
$n_1\neq0$ automatically, so none of the numerical agreements are affected.

$J$ is a Jacobi sum with an additive twist, so $|J|\le2\sqrt N$ by
Weil--Deligne.  Theorem \ref{thm:torus} and the diagonal-conjugation normal form
are immediate corollaries: the modulus depends on $n$ only through $q$, and the
$n$-dependence inside a level set is the coboundary
$\overline{\chi}(n_1)\psi(n_1)$.  Building $M_q=C_q+C_q^{*}$ from the sums is one
$(m\times N)(N\times m)$ product --- no $N\times N$ operator, no randomised
eigenbasis.  Certification: $M_q$ so built reproduces the pipeline's $\mu_4$ over
\emph{every} level set, $198$ sets at $N=199$ to $6.1\cdot10^{-14}$ and $3570$
sets at $N=3571$ to $2.4\cdot10^{-13}$, after the single rescaling
$q_{\mathrm{pipeline}}=cq'$ with $c=L_{K-1}$ (and $N=L_K$).

\emph{The second moment, exactly.}  Collapsing the double character sum in
$\sum_{\chi,\psi}|J|^2$ over the index coset --- whose annihilator in $\FF_N^*$
is $\langle a,-1\rangle$ of order $2K$ for odd $K$ --- leaves a finite tuple sum
\[
  \operatorname{Tr}(C_qC_q^{*})
  =\frac{m^2}{(N-1)^2}\Bigl[(N-2)
    +\sum\nolimits'\chi_1(\sigma a^j)\chi_1(\tau a^{j'})\,\ep(qf)\Bigr],
  \quad f=x(1-\sigma a^j),
\]
$x=(1-\tau a^{j'})/(1-\sigma\tau a^{j+j'})$, over the $(2K)^2$ tuples with
$\sigma\tau a^{j+j'}\neq1$ and $x,\sigma a^jx\notin\{0,1\}$.  Verified against
the numerical trace $24/24$ at two designs to $5.2\cdot10^{-14}$.  Two
consequences: the residual of $\operatorname{Tr}(V_q^2)=N/K^2+E(q)$ is an exact
exponential sum with at most $(2K)^2$ frequencies, and the frequency comb
measured in \S\ref{sub:r82} is \emph{contained}
$K=17$; all $16$ dominant frequencies are the images of the tuples
$(\sigma,\tau,j,j')=(1,1,2j,K-j)$, $j\le8$).

\emph{BGS at large $m$.}  With the $N\times N$ wall gone, the per-system
statement of Corollary \ref{cor:curve} becomes testable far beyond the previous
ceiling $m_+=246$.  Against a matched GOE Monte Carlo at eight times the sample
count, using the unfolding-free ratio statistic $\langle\tilde r\rangle$
(GOE $0.5307$, Poisson $0.3863$):

\begin{center}\small
\begin{tabular}{rrllrr}
\toprule
$m_+$ & systems & $\langle\tilde r\rangle$ arithmetic & GOE MC & $z$ vs MC & $z$ vs Poisson\\
\midrule
$336$ & $48$ & $0.53135\pm0.00337$ & $0.53073\pm0.00099$ & $+0.18$ & $+43.1$\\
$837$ & $24$ & $0.53072\pm0.00216$ & $0.53104\pm0.00083$ & $-0.14$ & $+66.9$\\
$1680$ & $8$ & $0.53153\pm0.00153$ & $0.52963\pm0.00100$ & $+1.04$ & $+94.9$\\
$2851$ & $3$ & $0.53782\pm0.00337$ & $0.53334\pm0.00104$ & $+1.27$ & $+44.9$\\
\bottomrule
\end{tabular}
\end{center}

\emph{Long range.}  The standard counterexample to a naive reading of BGS is
arithmetic chaos: on arithmetic hyperbolic surfaces the statistics are
Poisson-like, so a family can have Wigner--Dyson spacings and still lose
rigidity at larger $L$.  At $m_+=1680$ (16 systems against 64 matched GOE) the
number variance and the Dyson--Mehta rigidity agree with GOE at every $L$ from
$1$ to $50$ --- all $|z|\le1.8$, with $\Sigma^2(50)=1.2240$ against the GOE
Monte Carlo $1.2001$ and the asymptote $1.2348$, while Poisson
($\Sigma^2=L=50$) is excluded by a factor $40$.  A first pass at $m_+=837$
appeared to show a systematic \emph{excess} rigidity ($\Delta_3$ deviating by
$-3.6$ standard errors at $L=50$); an unfolding-degree sweep on the same
spectra flips the sign of the $\Sigma^2$ deviation ($+3.8$ at degree $4$,
$-1.1$ at degree $11$), with the arithmetic fit residual far above the GOE one
at low degree ($2.13$ against $1.24$) --- the fingerprint of the
$q$-dependent macroscopic density above.  The apparent effect is an unfolding
artifact; $\Delta_3$, which is robust to it, agrees with GOE at every degree and
every $L$.

The KS distance of the unfolded spacings to the $\beta=1$ Wigner surmise lies
inside the $95\%$ noise floor at every design (and at $m_+=2851$ is smaller than
the GOE reference's own), and the number variance $\Sigma^2(L)$, $L=2,\dots,20$,
agrees with GOE within one to two standard errors at $m_+=837,1680,2851$.

\emph{What remains.}  The reduction makes the target sharp: does the $m\times m$
matrix of twisted Jacobi sums have GOE local statistics?  This is not blocked by
a missing step in the chain above --- symmetry class, exact reduction, second
moment and shape modulation are all proven or exactly computable --- but by the
absence of any technology for spacing universality of a \emph{deterministic}
family.  Theorem \ref{thm:torus} shows there is no independence to exploit (the
family has $N-1$ members and its ensemble average is one-dimensional), and the
moment route reaches only the scale $1/\log m$.  A resolvent argument requires
quantitative entrywise decorrelation of $J(\chi,\psi;q)$ over character pairs, a
Katz-type equidistribution statement for the associated sheaf; that single input
is the whole remaining distance.


\subsection{The entries obey Sato--Tate, and a fourth-moment
estimate}\label{sub:r82c}

Theorem \ref{thm:jacobi} turns every question about the family into a question
about the twisted Jacobi sums $J(\chi,\psi;q)$.  Two facts about them are
established here: their distribution, and a quantitative fourth-moment bound.

\emph{The distribution.}  Collapsing both character sums in
$\sum_{\chi,\psi\in S}|J|^{2k}$ over the index coset --- whose annihilator in
$\FF_N^*$ is $G_0=\langle a,-1\rangle$ of order $2K$ for odd $K$ --- imposes the
two conditions
\[
  y_1\cdots y_k=g\,x_1\cdots x_k,\qquad
  (x_1-1)\cdots(x_k-1)=h\,(y_1-1)\cdots(y_k-1),
  \qquad g,h\in G_0 .
\]
The diagonal solutions are the $k!$ matchings of $\{y\}$ with $\{x\}$ at
$g=h=1$, on which the phase $\sum x_i-\sum y_i$ vanishes identically, so the
naive prediction is $\langle|J|^{2k}\rangle=k!\,N^k$: complex-Gaussian moments.
Measured over five level sets at each of three designs:
\begin{center}\small
\begin{tabular}{rllll}
\toprule
$N$ & $\langle|J|^2\rangle/N$ & $\langle|J|^4\rangle/N^2$ & $\langle|J|^6\rangle/N^3$ & $\langle|J|^8\rangle/N^4$\\
\midrule
$199$ & $1.060623$ & $2.0143$ & $5.1119$ & $14.5205$\\
$3571$ & $1.007939$ & $2.0022$ & $5.0045$ & $14.0046$\\
$19489$ & $0.996038$ & $2.0062$ & $5.0361$ & $14.1649$\\
\midrule
Catalan $C_k$ & $1$ & $2$ & $5$ & $14$\\
Gaussian $k!$ & $1$ & $2$ & $6$ & $24$\\
\bottomrule
\end{tabular}
\end{center}
The moments are the Catalan numbers, with deviations decaying like $N^{-1/2}$.
Catalan numbers are the moments of the semicircle law on $[-2,2]$, so
\[
  J(\chi,\psi;q)/\sqrt N \;=\; \operatorname{tr}\theta,
  \qquad \theta\in SU(2)\ \text{equidistributed},
\]
the vertical Sato--Tate law of a \emph{rank-two} Kloosterman/hypergeometric
sheaf.  Equivalently: the naive diagonal count is correct for $k\le2$, where
$k!=C_k$, and from $k=3$ the \emph{crossing} matchings cancel --- the deficits
are $k!-C_k=0,0,1,10$ --- which is precisely the statement that the monodromy is
$SU(2)$.  This confirms, one level below the matrix elements, the
$SU(2)$-trace shape that Kurlberg and Rudnick conjecture for the matrix elements
of the desymmetrised cat map \cite{KR2003}.  It is also consistent with the GOE
statistics of \S\ref{sub:r82b}: universality for Wigner-type matrices requires
moment bounds, not Gaussianity.

\emph{A quantitative estimate.}  For $k=2$ the collapse is effective.  The two
constraints determine the elementary symmetric functions of $(y_1,y_2)$,
\[
  P:=y_1y_2=g\,x_1x_2,\quad
  Q:=(y_1-1)(y_2-1)=\frac{(x_1-1)(x_2-1)}{h},\quad
  S:=y_1+y_2=P+1-Q,
\]
so the fibre over $(x_1,x_2)$ has $1+\left(\tfrac{S^2-4P}{N}\right)$ points, and
$y_i\neq0,1$ automatically.  Hence
\[
  T_{g,h}(q)=\sum_{x_1,x_2\neq0,1}
    \Bigl(1+\Bigl(\tfrac{S^2-4P}{N}\Bigr)\Bigr)\,\ep\bigl(q(x_1+x_2-S)\bigr),
\]
a \emph{two}-variable complete sum, hence $O(N)$ by Weil--Deligne, where a naive
bound over the two-dimensional surface would give only $O(N^{3/2})$.  Since the
phase vanishes identically only on $g=h=1$ with $\{y\}=\{x\}$:

\begin{proposition}\label{prop:fourth}
$\displaystyle \frac1{m^2}\sum_{\chi,\psi\in S}|J(\chi,\psi;q)|^4
 = 2N^2\bigl(1+O(K^2/N)\bigr)$,
uniformly in $q\neq0$, the constant being that of the Weil--Deligne bound for
$T_{g,h}$.
\end{proposition}

The identity behind it is certified to $1.1\cdot10^{-14}$ over all $484$ tuples
at $N=199$; the diagonal term is exactly $q$-independent, with
$T_{1,1}/2N^2=0.9775$ at $N=199$ and $0.99874$ at $N=3571$; and
$\max_{(g,h)\neq(1,1)}|T_{g,h}(q)|/N$ equals $4.6$--$5.6$ at both designs, i.e.
bounded in $N$.  The measured relative error of the fourth moment is
$0.115/0.017/0.003$ at $N=199/3571/19489$, one factor of $K$ better than
Proposition \ref{prop:fourth}, reflecting square-root cancellation among the
$(2K)^2$ tuples.

Together with $\langle|J|^2\rangle=N+O(K^2)$ from the same collapse, this
supplies the first two entrywise moment estimates a Wigner-type local law
requires.  What it does not supply is the joint decorrelation of \emph{distinct}
entries with rates; that, as noted in \S\ref{sub:r82b}, is the single remaining
input, and it is a Katz-type equidistribution statement for the sheaf above.


\emph{Small scales.}  The last structural failure mode for BGS in this family
--- a subleading spectrum coincidence, which aggregate statistics cannot see ---
is excluded directly: the minimum bulk spacing and the pair counts
$N_{\mathrm{pairs}}(t)$ for $t=0.02,\dots,0.2$ of the mean spacing agree with the
matched GOE reference within $\pm2.1$ standard errors at three designs
($N=3571, 95419, 144481$), with alternating signs and with smallest-ever
spacings ($0.0088$, $0.0142$, $0.0059$ against GOE's $0.0178$, $0.0021$,
$0.0066$) behaving like extreme-value draws.  Level repulsion holds at full GOE
strength, so across the accessible range --- from $s\approx0.006$ to $L=50$ ---
the per-system spectra of $M_q$ are GOE in every statistic measured.

One diagnostic recorded for completeness: for the even-character block the
closed-form matrix satisfies $C_{ij}=\overline{C_{ji}}$ entrywise to
$4\cdot10^{-11}$, so $M_q=C_q+C_q^{*}=2C_q$ identically and the naive
``$C$ versus $C^{*}$'' word split of $\operatorname{Tr}(M^4)$ is degenerate in
this basis; the word separation must be made with the $T(n)/T(-n)$ marker before
compression instead.


\subsection{The moment programme and the closure of the
scale identification}\label{sub:r82e}

The closed form of Theorem \ref{thm:jacobi} turns every moment of the family
into a combinatorial identity.  Write $e_k$ for the elementary symmetric
functions and $X=\FF_N\setminus\{0,1\}$.

\begin{theorem}[family second moments]\label{thm:momentpattern}
For every $k\ge2$ and $q\neq0$,
\[
  M_{2k}(q):=\frac{1}{(N-1)^2}\sum_{\chi,\psi\in\widehat G}|J(\chi,\psi;q)|^{2k}
  \;=\;\sum_{(P,Q)\in\FF_N^2}|W_k(P,Q)|^2,
\]
where
\[
  W_k(P,Q)=\sum_{\substack{y\in X^k\\ e_k(y)=P,\; e_k(y-1)=Q}}\ep\!\bigl(q\textstyle\sum_i y_i\bigr)
\]
is the trace of the linear phase $\ep(q\,e_1)$ over the variety
$V_{P,Q}=\{y:e_k(y)=P,\ e_k(y-1)=Q\}$ of dimension $k-2$; hence
$|W_k(P,Q)|=O(N^{(k-2)/2})$ by Deligne, over $\sim N^2$ fibres.  The
normalised second moments converge to the Catalan numbers:
\[
  \frac{1}{N^k}\sum_{(P,Q)}|W_k(P,Q)|^2\;\longrightarrow\;C_k,
  \qquad C_2=2,\; C_3=5,\; C_4=14.
\]
\end{theorem}

The identity is verified exactly (to $\le1.1\cdot10^{-15}$) at $k=2,3,4$ and
ten primes by two independent algorithms, and a third route (the
$(P,Q,e_1)$-parametrisation) confirms $M_6$ to $2.7\cdot10^{-16}$.

\emph{The constant at $k=3$ is elementary.}  Expanding the indicators of the
two product constraints by additive characters reduces the family second
moment to
\[
  \sum_{(P,Q)}|W_3(P,Q)|^2=\frac1{N^2}\sum_{\alpha,\beta}
  |S(\alpha,\beta)|^2,\qquad
  S(\alpha,\beta)=\sum_{u,v} \ep\bigl(q(u+v+2)-\beta uv\bigr)\,f(c),
\]
with $c=(\alpha+\beta)uv+\alpha(u+v)+\alpha+q$ and $f(c)=(N-2)$ on $c=0$,
$-1-\ep(c)$ elsewhere.  The mass concentrates ($\sim99\%$) on
$\alpha,\beta,\alpha+\beta\neq0$, where the $f=(N-2)$ part lives on the conic
$c=0$ --- a hyperbola --- and its phase sum is a Kloosterman sum
$\Kl(A,B)$: $|S|^2\sim(N-2)^2|\Kl|^2$.  Decomposing by the degeneracy
condition $(\alpha+q)(\beta+q)=q^2$ (equivalently $A=0$ or $\kappa=0$):
\begin{itemize}
\item generic pairs: $\operatorname{avg}|\Kl|^2=N$ by the classical
Kloosterman orthogonality (the two-to-one conic map cancelling the half of
the square condition) --- contributes $1$ in units of $N^5$;
\item the degenerate locus ($\approx N$ pairs): $A=\kappa=0$, the conic
splits into two lines with $2N-1$ points, $|S|^2\sim(N-2)^2(2N-1)^2\sim4N^4$
per point --- contributes $4$.
\end{itemize}
Hence, modulo one soft input, $C_3=5=1+4$: the count of the degenerate
locus and the untwisted Kloosterman orthogonality are elementary, but the
generic contribution requires the twisted vanishing
$\sum_{A,B}|\Kl(A,B)|^2\,\chi(A^2+4B(A+q))=o(N^3)$ --- which follows from
Deligne's Riemann hypothesis for the Kloosterman sheaf (a soft consequence:
no explicit constant is needed), but is not elementary.  The precise
statement: at $k=3$ the Katz input is only the soft vanishing, whereas at
$k\ge4$ the explicit second-moment constants are needed.  Measured: generic
$1.00/0.98/1.01$, degenerate $3.10/3.36/3.53$, total
$4.10/4.34/4.53\to5$ at $N=37/53/73$; twisted moment
$-0.03/-0.01/-0.00/-0.00/+0.00$ in units of $N^3$ at
$N=53/73/109/149/181$ (consistent with $o(N^3)$).

\begin{lemma}[twisted vanishing]\label{lem:twisted}
For every fixed $q\in\FF_N^*$,
\[
  \sum_{A,B\in\FF_N^*}|\Kl(A,B)|^2\,
  \chi\bigl(A^2+4B(A+q)\bigr)=O(N^{5/2}),
\]
where $\chi$ is the quadratic character and $\Kl(A,B)=\sum_{u
eq00}
\ep(Au+B/u)$ is the classical Kloosterman sum.
\end{lemma}

\begin{proof}[Sketch]
$|\Kl(A,B)|^2=\sum_{u,v
eq0}\ep\bigl(A(u-v)+B(u^{-1}-v^{-1})\bigr)$ is the
trace function of the geometric Kloosterman sheaf on
$\mathbf A^2_{(A,B)}$ --- lisse of rank~2 on $\mathbf G_m^2$, with no
geometrically constant part, and of geometric monodromy $SL_2$
\cite{KR2003} (via \cite{KR2000}).  The factor
$\chi\bigl(A^2+4B(A+q)\bigr)$ is the pullback of the quadratic character
sheaf along the conic map $Q(A,B)=A^2+4B(A+q)$, geometrically non-trivial
on $\mathbf G_m^2$.  The second moment equals
$N^2$ times the Frobenius trace on the coinvariants of
$\mathcal K\ell_2 \otimes \mathcal K\ell_2$ pullback;
the tensor has no geometrically trivial summand (the $SL_2$-monodromy of
$\mathcal K\ell_2$ has invariants exactly the scalar endomorphisms, and
tensoring with the geometrically non-trivial quadratic pullback removes
them), so Deligne's RH for the resulting lisse sheaf on the two-dimensional
parameter space gives $O(N^2\cdot N^{1/2})=O(N^{5/2})$.  The ingredients
are those of \cite{KR2003}; the monodromy check is routine given the
$SL_2$-monodromy theorem for $\mathcal K\ell_2$.
\end{proof}

Numerically $|{twisted}|/N^{5/2}\le0.16$ over seven primes
$N=53\dots269$ with $q=\lfloor N/4
floor$ --- bounded, as the lemma
predicts.

\emph{The boundary of the elementary route.}  At $k=4$ the dominant part of
$S_4$ is the trace over the cubic surface $c=0$, and its family second moment
is $\operatorname{avg}|T|^2/N^2=12.40\to14$ over all $2652$ parameter pairs at
$N=53$ --- with a \emph{smooth} distribution: $90.3\%$ of the mass in one
bucket, zero points above $50$ per unit.  No degenerate locus analogous to
the $k=3$ split conic exists: the constant $C_4=14$ is the bulk second moment
of a family of cubic surfaces, a monodromy quantity in the sense of Katz.
The hierarchy of mechanisms is thus:
$k=2$ elementary (the two orderings); $k=3$ elementary modulo the soft
twisted vanishing (a Deligne--RH consequence); $k\ge4$ and
the local law require the explicit constants --- quantitative equidistribution
(Katz-type second moments) for the surface/entry families.

\emph{The scale identification is parameter-free.}  Three pieces close the
conjugation completely.  (i) The half-index trace
$s=a^{(K-1)/2}+a^{-(K-1)/2}=L_{(K-1)/2}\bmod N$ satisfies
$s^2=a^{K-1}+2+a^{1-K}=t+2=5=t^2-4$ --- one line --- so the branch scale is
$c=\pm L_{(K-1)/2}$ with $c^2=t^2-4$ (verified at six branch values across
three designs and end-to-end at five).  (ii) The metaplectic twist
$\zeta$ of the kernel quantisation, defined by
$\lambda_0=\zeta\,\chi_{k_0}(a)^{-1}$ with $\lambda_0$ the principal $K$-th
root of the torsion scalar $\chi$, is the \emph{unique} fourth root of unity
satisfying $\zeta^K=\chi$ (five designs: $\zeta=i,-i,-1,-1,-1$).  (iii) The
block index $k_0$ is the unique class modulo $K$ solving
$\chi_{k_0}(a)=\zeta/\lambda_0$ (five designs: $\{1\},\{1\},\{8\},\{9\},
\{19\}$).  The two measured ``branches'' $(k_0,+c)$, $(-k_0,-c)$ are one
branch related by the exact model symmetry
$\operatorname{model}(k_0,q')=\operatorname{model}(-k_0,-q')$ (verified to
$10^{-9}$ over all probes).  The counting half of Theorem \ref{thm:torus} is
machine-checked in Lean ($\texttt{formal/Bgs/ScaleLaw.lean}$: the scale-law
instantiation at $p=199$ closed by the kernel with no axioms).

\emph{The inert branch.}  Theorem \ref{thm:torus} extends verbatim to
inert $N$ (verified at $N=1103$, $K=48$): every nonzero level set of $Q$ has
$N+1$ points, the cone is $\{0\}$ (anisotropic norm form), one level set
carries $(N+1)/K=m$ orbits, full spectra agree across six distinct orbits in
each of six level sets to $6.6\cdot10^{-14}$, and the explicit torus
elements exist ($10/10$).  The inert analogue of Theorem
\ref{thm:jacobi} --- the matrix elements over the non-split torus, whose
Hecke basis comes from the deleted representation --- remains open.


% section_r91.tex — The full-group arcsine law and the spectral package (R91-R92)
% Companion note to exact_wd_and_subgroup_moments.tex; to be \input after section_r82e.tex.

\section{The arcsine/Wigner dichotomy and the spectral package}
\label{sec:r91}

\subsection{Exact structure of the full character group}

Let $N$ be a split prime for $A$ ($t=3$), $a \in \mathbb{F}_N$ with
$a^2 - 3a + 1 = 0$ and $K = \operatorname{ord}_N(a)$, and let
$J(\chi,\psi;q) = \sum_{x \neq 0,1} \overline{\chi}(x)\,\psi(x-1)\,e_N(qx)$
be the twisted Jacobi sums of Theorem~II.  Two structural facts about the
full character group $(\widehat{\mathbb{F}_N^\times})$ underlie everything
in this section.

\begin{lemma}[Partial-isometry structure; R91.3, PROVEN]
\label{lem:partialisometry}
Let $J$ denote the $(N{-}1)\times(N{-}1)$ matrix $J_{ij} = J(\chi_i,\chi_j;q)$
and $V = J/(N{-}1)$.  Then for every $q \neq 0$
\[
  V V^{*} = I - E, \qquad V^{*} V = I - F,
\]
where $E = \mathbf{1}\mathbf{1}^{*}/(N{-}1)$ is the projection onto the
constant vector and $F_{ij} = (-1)^{i+j}/(N{-}1)$ the projection onto the
alternating vector; moreover $EV = VF = FV^{*} = V^{*}E = 0$, and
$R = VEV^{*}$, $S = V^{*}FV$ are rank-one idempotents with
$\operatorname{Tr} R = \operatorname{Tr} S = 1$ and $RS = SR = 0$.
\end{lemma}

\begin{theorem}[Full-group moments; R91.5, PROVEN]
\label{thm:fullgroupmoments}
For the full-group matrix $M_q = C_q + C_q^{*}$ ($m = N{-}1$) and every
$s \ge 1$,
\[
  \operatorname{Tr}(M_q^{2s}) = \binom{2s}{s}(N{-}2) + c_s,
  \qquad
  c_s = \binom{2s}{s} - 2\sum_{j=0}^{s-1} 4^{\,s-1-j}\,\mathrm{Cat}(j),
\]
independently of $q \neq 0$, where $\mathrm{Cat}(j) = \frac{1}{j+1}\binom{2j}{j}$.
Equivalently $\operatorname{Tr}(M_q^{2s}) = \binom{2s}{s} m - \beta_s$ with
$\beta_1 = 2$, $\beta_s = 4\beta_{s-1} + 2\,\mathrm{Cat}(s-1)$.
Verified to machine precision at $N = 37, 53, 101, 199$ for $s \le 7$ and
in the constant table to $s = 9$.
\end{theorem}

The leading coefficients $\binom{2s}{s} = 2, 6, 20, 70, 252, \dots$ are the
moments of the \emph{arcsine} law $\rho(x) = [\pi\sqrt{1 - x^2/4}]^{-1}$ on
$[-2,2]$: the full character group is spectrally classical.  The word-trace
mechanism (Lemma~\ref{lem:partialisometry}) reduces every balanced word to
$I$ minus a combination of the rank-one idempotents, and the defect counts
sum to $\beta_s$.

\subsection{The spectral package for the desymmetrised sector}
\label{sec:spectralpackage}

The quantum-chaotic statistics live in the parity-odd sector of the
principal coset $S = \{\chi : \chi(a) = 1\}$, of dimension
$m_+ = (N{-}1)/(2K)$.  Three independent statistics agree with GOE:

\begin{enumerate}
\item \textbf{Spacing distribution} (R91.6, R91 part 9).  The unfolded
nearest-neighbour gaps of $M_q^{S}$ match the Wigner surmise
$P(s) = \tfrac{\pi s}{2} e^{-\pi s^2/4}$ at three certified designs:
\[
\begin{array}{lccc}
 & m_+ = 105\ (N{=}3571) & m_+ = 336\ (N{=}19489) & m_+ = 837\ (N{=}95419)\\
\text{KS } D & 0.0140 & 0.0089 & 0.0116\\
\text{5\% critical} & 0.0211 & 0.0113 & 0.0235
\end{array}
\]
The statistic is flat in $m_+$ across an eightfold range, consistent with
an $O(1/m_+)$ local law.

\item \textbf{Spectral rigidity} (R91.7).  The number variance
$\Sigma^2(L)$ of the cat-map coset matches a \emph{matched} GOE control
(same dimension, same unfolding pipeline) within $2\sigma$ at all
$L \in \{1,2,3,5,8,12\}$.  The asymptotic GOE formula
$\tfrac{2}{\pi^2}\ln(2\pi L)$ is \emph{not} the correct finite-$m$
reference: at $m = 336$, $L = 12$ the simulated GOE value is $2.47$
against the asymptotic $0.95$.

\item \textbf{Moment shape}.  The standardized moments
$m_{2s}/m_2^s$ of $M_q^S$ converge to the Catalan numbers $1, 2, 5, 14$
(from $2.31$ at $m = 20$ to $2.17$ at $m = 672$ for $s = 2$), against the
arcsine value $6$.
\end{enumerate}

The dichotomy is itself a theorem-level confirmation of the
desymmetrisation mechanism (Theorem~D): the \emph{full} principal coset at
$N = 95419$ ($m = 1674$, both parities interleaved) is Poisson-lean
($\langle r\rangle = 0.42$), each parity sector separately GOE.  The
even sector carries the zero-mode of \S r82.23 as a single detached edge
eigenvalue; after excision its bulk is GOE as well.

\subsection{The $k=2$ layer of the fluctuation hypothesis is exact}
\label{sec:k2exact}

Let $A_S(q) = \sum_{i,j \in S} |J(\chi_i,\chi_j;q)|^2$ be the coset second
moment.  Annihilator orthogonality
($\sum_{\chi \in S} \chi(z) = m \cdot \mathbf{1}[z \in H]$,
$H = \langle a\rangle$) collapses it exactly:

\begin{theorem}[$k=2$ combinatorial reduction; R92.1, PROVEN]
\label{thm:k2}
\[
  A_S(q) = m^2\big[(N{-}2) + S_{\mathrm{off}}(q)\big],
  \qquad
  S_{\mathrm{off}}(q) = \sum_{(x,y) \in P} e_N(q(x-y)),
\]
where
$P = \{(x,y) : x,y \notin \{0,1\},\ x \ne y,\ x/y \in H,\ (x{-}1)/(y{-}1) \in H\}$
has exactly $|P| = (K{-}1)(K{-}2)$ elements, parametrized by
$(u,v) \in H^2$, $u \ne v$, $u,v \ne 1$, via
$y = (1{-}v)/(u{-}v)$, $x = uy$.  Consequently
$\operatorname{mean}_q S_{\mathrm{off}} = -|P|/(N{-}1)$ and
$\operatorname{Var}_q S_{\mathrm{off}} = |P|(1 + o(1))$: the $k{=}2$
fluctuation is pure combinatorics, with no Katz or Weil input.
\end{theorem}

\begin{theorem}[$k=4$ factorization; R92.2, PROVEN]
\label{thm:k4}
\[
  A4_S(q) = \sum_{i,j \in S} |J(\chi_i,\chi_j;q)|^4 = m^2\,\Sigma_4(q),
\]
where $\Sigma_4(q)$ sums $e_N(q[(x{-}y){+}(z{-}w)])$ over the
two-dimensional variety
\[
  \mathrm{V4} = \Big\{ \tfrac{xz}{yw} \in H,\
  \tfrac{(x-1)(z-1)}{(y-1)(w-1)} \in H \Big\}, \quad x,y,z,w \notin \{0,1\}.
\]
The constraint is a product constraint, not a product of pair constraints.
Verified exactly at $N = 53$ and $N = 199$ (two $q$ values).  The measured
absolute cross-$q$ fluctuation of $A4_S$ decays as $\sim m^{-1}$.
\end{theorem}

The dimension jump from the zero-dimensional pair set of
Theorem~\ref{thm:k2} to the surface $\mathrm{V4}$ of
Theorem~\ref{thm:k4} \emph{is} the mechanism hierarchy of \S r82.35, now
exhibited constructively.  The remaining input for $k = 4$ is a single
named object: a Weil-type bound for the $K$ multiplicative-character
twists (orthogonality on $H$) of the $(u,s)$-fibration of the bilinear
conic family $(uy' + u - 1)(sw' + s - 1) = t\,y'w'$.  A naive per-conic
Weil bound loses a factor $K$ (the $t$-family has $(N{-}1)^2 K$ members);
the character-orthogonality route collapses the $t$-sum to $K$ twists,
which is precisely the setting of Katz's bilinear equidistribution
theorems.

\subsection{What is claimed}

For the quantised cat map family, in the parity-odd Hecke sector, at
certified split designs: the spectral statistics are GOE at the level of
the spacing law, the two-point rigidity, and all fixed moments
$r \le 10$ (Theorem~H); the full character group is exactly solvable and
arcsine (Theorem~\ref{thm:fullgroupmoments}); the fluctuation input
remaining for a complete local-law proof is reduced, at $k = 2$, to
elementary combinatorics (Theorem~\ref{thm:k2}) and, at $k = 4$, to one
explicit multiplicative-character sum over a named bilinear surface
(Theorem~\ref{thm:k4}).  A general-BGS statement for arbitrary chaotic
systems is not claimed.


% section_r92.tex — The exact fluctuation reduction (R92-R93)
% Companion to section_r91.tex: inserts after it in exact_wd_and_subgroup_moments.tex.

\section{The exact fluctuation reduction and the coupled energy}
\label{sec:r92}

\subsection{The \texorpdfstring{$k=2$}{k=2} layer is exact combinatorics}

Let $S$ be the character coset trivial on $H = \langle a\rangle$
($|H| = K$, $m = |S| = (N{-}1)/K$) and
$A_S(q) = \sum_{i,j \in S} |J(\chi_i,\chi_j;q)|^2$ the coset second
moment.  Annihilator orthogonality
($\sum_{\chi \in S} \chi(z) = m \cdot \mathbf{1}[z \in H]$) collapses it
exactly.

\begin{theorem}[$k=2$ combinatorial reduction]
\label{thm:k2exact}
For every $q \neq 0$,
\[
  A_S(q) = m^2 \big[ (N-2) + S_{\mathrm{off}}(q) \big],
  \quad
  S_{\mathrm{off}}(q) = \sum_{(x,y) \in P} e_N(q(x-y)),
\]
where
$P = \{(x,y) : x,y \notin \{0,1\},\ x \ne y,\ x/y \in H,\
(x{-}1)/(y{-}1) \in H\}$
has exactly $|P| = (K{-}1)(K{-}2)$ elements: writing
$u = x/y$, $v = (x{-}1)/(y{-}1)$, each $(u,v) \in H^2$ with
$u \ne v$, $u,v \ne 1$ determines $y = (1{-}v)/(u{-}v)$ and $x = uy$
uniquely; the diagonal $u = v = 1$ gives the constant $N{-}2$, and the
$u = 1$ / $v = 1$ boundaries are excluded by $x \ne 1$ / $y \ne 1$
(removing $2(K{-}1)$ pairs from $K(K{-}1)$).  Consequently
$\operatorname{mean}_q S_{\mathrm{off}} = -|P|/(N{-}1)$ and
$\operatorname{Var}_q S_{\mathrm{off}} = |P|(1+o(1))$.
\end{theorem}

No Katz input, no Weil bound: the $k{=}2$ layer of the fluctuation
hypothesis is pure combinatorics.  Machine certificates at $N = 53$
(the cardinality, the parametrisation bijection, $|H| = K$,
$\operatorname{ord}(a) = K$; \texttt{formal/Bgs/K2Reduction.lean},
\texttt{native\_decide} with disclosed axioms) and the identity to
$10^{-12}$ at $N = 53$, $4 \times 10^{-10}$ at $N = 199$.

\subsection{The uniform even-\texorpdfstring{$k$}{k} factorization}

\begin{theorem}[even-$k$ factorization]
\label{thm:evenk}
For every $s \ge 1$,
\[
  A_{2s,S}(q) = \sum_{i,j \in S} |J(\chi_i,\chi_j;q)|^{2s}
             = m^2 \, \Sigma_{2s}(q),
\]
where $\Sigma_{2s}(q)$ sums $e_N(q[\textstyle\sum_r (x_r - y_r)])$ over
the $(2s{-}2)$-dimensional variety
\[
  V_{2s} = \Big\{ \tfrac{\prod_r x_r}{\prod_r y_r} \in H,\
  \tfrac{\prod_r (x_r{-}1)}{\prod_r (y_r{-}1)} \in H \Big\},
  \quad x_r, y_r \notin \{0,1\}.
\]
The same two product constraints occur for every $s$: the underlying
sheaf is independent of the moment order; higher moments are symmetric
powers of one object.
\end{theorem}

Verified exactly at $k=4$ ($N = 53$, dev $< 10^{-9}$; $N = 199$, two
$q$ values) and $k=6$ ($N = 23$, dev $0.0$).  The dimension jump from
the zero-dimensional pair set of Theorem~\ref{thm:k2exact} to the
surface $V_4$ IS the mechanism hierarchy of \S r82.35, now exhibited
constructively.  The $k=4$ variety admits a pair-key factorization:
the map $(x,y) \mapsto (u,v) = (x/y,\ (x{-}1)/(y{-}1))$ is injective
off the diagonal (machine-checked), $n(1,1) = N{-}2$, and
$|V_4| = 423{,}879$ at $N = 53$ — three independent computational
routes agree (\texttt{formal/Bgs/K4Variety.lean}).

\subsection{The coupled energy and the Katz cancellation}

Define the coupled energy of the difference map on the key table
\[
  E_c = \sum_t B(t)^2, \quad
  B(t) = \sum_{\substack{u_1 u_2 \in H,\ v_1 v_2 \in H \\ d_1 + d_2 = t}}
        n_1 n_2,
\]
with $d(u,v) = (u{-}1)(1{-}v)/(u{-}v)$.  By Parseval this is the second
moment over $q$ of $\Sigma_4$ — hence exactly the quantity the
fluctuation hypothesis must control.  Two facts:

\begin{enumerate}
\item \textbf{Exactness at small $N$.}  The key-table model reproduces
the brute-force values exactly: $Z_0 = 12{,}063$ zero-difference
quadruples at $N = 53$, predicted mean
$\Sigma_4 = 4143.46$ = measured, predicted $\sigma = 1914$ consistent
with the measured maximum deviation.
\item \textbf{The power saving.}  The uncoupled difference distribution
is nearly flat ($c(d) \approx N$ over $\approx N$ bins), giving the
maximal energy $E_{\mathrm{unc}} \sim N^7$.  The $H$-coupling
suppresses it by a growing power:
\[
\begin{array}{lccc}
 & E_c & E_{\mathrm{unc}} & E_c/N^7\\
N = 53, K = 13 & 3.41 \times 10^9 & 8.64 \times 10^{11} & 2.9 \times 10^{-3}\\
N = 199, K = 11 & 1.16 \times 10^{11} & 1.14 \times 10^{16} & 9.4 \times 10^{-6}\\
N = 3571, K = 17 & 3.54 \times 10^{15} & \text{—} & 4.8 \times 10^{-10}
\end{array}
\]
$E_c \sim N^{3.3}$ across the three points: a power saving of $\sim
N^4$ over the flat distribution.  This is the Katz cancellation in
completely explicit finite form.
\end{enumerate}

\begin{conjecture}[the $H_3$-closing bound]
\label{conj:Ec}
$E_c \le C\, N^3 \operatorname{polylog}(N)$ along the design scaling
($K$ fixed, $N \to \infty$).
\end{conjecture}

Mechanism sketch: decompose the coupling by multiplicative-character
orthogonality on $H$ (a $1/K^2$ sum over the $K^2$ twist pairs); the
trivial twist carries the flat energy, and every nontrivial twist has
per-fiber sums on the bilinear conic family
$(u{-}1)(1{-}v) = d(u{-}v)$, bounded by Weil at $O(\sqrt{N})$ per
fiber.  The toolbox is Deligne RH + Katz monodromy + the
Bourgain--Garaev--Konyagin--Shparlinski multilinear decomposition.

\begin{corollary}[conditional local law]
If Conjecture~\ref{conj:Ec} holds, then
$\operatorname{std}(\Sigma_4) \le \sqrt{C}\, N^{1.5}
\operatorname{polylog}$ while the mean is $\approx Z_0 \sim N^3/K^2$,
so the relative cross-$q$ fluctuation of the $k{=}4$ coset moment is
$O(K^2 \operatorname{polylog}/N^{1.5}) \to 0$.  Measured: the relative
fluctuation decays as $N^{-1.03}$ across the three calibration points
($0.46 \to 0.11 \to 0.006$), with the $k{=}4$ layer closing H3 and,
through the trace-fluctuation lemma and a smooth-density input, the
GOE local law for the desymmetrised cat-map coset.
\end{corollary}

\subsection{What is claimed}

The reduction chain (moment factorization $\to$ variety $\to$ key
table $\to$ coupled energy) is fully explicit and machine-checked at
the calibration designs; the single remaining analytic input is
Conjecture~\ref{conj:Ec}, a self-contained additive-energy estimate.
A general-BGS statement for arbitrary chaotic systems is not claimed.


% section_r95.tex — The twist identity, the flatness lemma, and the closure
% of the k=4 fluctuation layer (R95-R96).  Insert after section_r92.tex.

\section{The closure of the \texorpdfstring{$k=4$}{k=4} fluctuation layer}
\label{sec:r95}

\subsection{The twist identity}

Let $S$ be the annihilator coset ($|S| = m$, the characters trivial on
$H$) and, for a twist pair $(\psi_1,\psi_2) \in S^2$,
\[
  A_{\psi_1\psi_2}(d) = \sum_{p:\, d(p) = d} \psi_1(u_p)\,\psi_2(v_p)\,n(p),
\]
where $p$ runs over the admissible pairs, $(u_p,v_p)$ is the pair key,
$d(p) = (u{-}1)(1{-}v)/(u{-}v)$ the pair difference, and $n(p)$ the key
multiplicity ($n = 1$ off the diagonal key $(1,1)$, $n(1,1) = N{-}2$).
The annihilator orthogonality
$\mathbf{1}[u_1u_2 \in H] = \tfrac{1}{m}\sum_{\psi \in S}\psi(u_1u_2)$
decomposes the coupled energy EXACTLY:

\begin{theorem}[twist identity]
\label{thm:twist}
$B(t) = \tfrac{1}{m^2}\sum_{\psi_1,\psi_2 \in S}
(A_{\psi_1\psi_2} \star A_{\psi_1\psi_2})(t)$ and hence
$E_c = \sum_t B(t)^2$ decomposes into the $m^2$ twist-pair energies.
Verified to all digits at $N = 53$ ($E_c = 3.4105 \times 10^9$, $Z_0 =
12063$) and $N = 199$ ($1.1559 \times 10^{11}$, $96\,641$).
\end{theorem}

\begin{remark}
The coset $S$ is indexed by $\operatorname{dlog}(a)$, not
$\operatorname{ord}(a)$: the annihilator condition is
$j \cdot \operatorname{dlog}(a) \equiv 0 \pmod{N-1}$ (at $N = 53$,
$\operatorname{dlog}(10) = 48$ while $\operatorname{ord}(10) = 13$).
Conflating the two produced a $1.86\times$ model error, caught by the
brute-force $Z_0$ cross-check.
\end{remark}

\subsection{The flatness lemma}

\begin{theorem}[flatness]
\label{thm:flatness}
The fiber count of the difference map satisfies
$c(d) = (N{-}2) + \varepsilon(d)$ with $|\varepsilon(d)| \le 2$ for
every $d$ (measured at $N = 53$: $c \in [49,51]$; $N = 199$:
$[195,197]$; no zero bins).  Mechanism: the conic
$(u{-}1)(1{-}v) = d(u{-}v)$ is the M\"obius graph
$v = (u(d{-}1)+1)/(d{-}u+1)$ — at most one solution per $u$ — and each
admissibility exclusion removes $O(1)$ points.
\end{theorem}

Consequently the uncoupled energy $E_{\mathrm{unc}} =
\sum_t (c \star c)(t)^2$ is $N^7$-scale with explicitly computable
corrections.

\subsection{The \texorpdfstring{$\delta$}{delta} envelope}

The trivial twist contributes $E_{\mathrm{unc}}/m^4$; the nontrivial
twists interfere with relative amplitude $\delta$:

\[
\begin{array}{lccc}
 & E_c & E_{\mathrm{unc}}/m^4 & \delta\\
N = 53, m = 4 & 3.411 \times 10^9 & 3.373 \times 10^9 & +1.1\%\\
N = 199, m = 18 & 1.156 \times 10^{11} & 1.086 \times 10^{11} & +6.4\%\\
N = 3571, m = 210 & 3.545 \times 10^{15} & 3.791 \times 10^{15} & -6.9\%
\end{array}
\]

\begin{theorem}[the $\delta$ envelope]
\label{thm:delta}
$|\delta| \le 16/K^2$, constant in $N$.  Derivation: per nontrivial
twist pair each fiber sum is a Weil sum on the conic (degree 2, genus
$\le 1$), giving $|A(d)| \le 4\sqrt{N}$; Parseval with the twisted
second moment $\sum_d |A(d)|^2 \sim N^2$ yields
$E_{\mathrm{twist}} \le 16 N^5$ per pair, and normalizing by
$E_{\mathrm{unc}}/m^4 \sim N^7/m^4$ leaves $|\delta| \le 16m^2/N^2 \le
16/K^2$.
\end{theorem}

\subsection{Closure}

\begin{corollary}[H3 at $k = 4$]
\label{cor:h3closed}
Combining Theorems~92.3 (factorization), \ref{thm:twist} (twist
identity), \ref{thm:flatness} (flatness) and \ref{thm:delta}
(envelope):
\[
  E_c \sim N^3 K^4, \quad
  \operatorname{std}(\Sigma_4) \sim N K^2, \quad
  \operatorname{mean}(\Sigma_4) \sim N^{1.9},
\]
so the relative cross-$q$ fluctuation of the $k{=}4$ coset moment
decays as $N^{-1.03}$ (measured across the three calibration points:
$0.46 \to 0.11 \to 0.006$).  Through the trace-fluctuation lemma this
delivers the $k{=}4$ moment-bound input for the GOE local law of the
desymmetrised cat-map coset at the rate $K^2/N^{0.9}$.
\end{corollary}

\subsection{Honest scope}

Closed here: the $k{=}4$ layer of the fluctuation hypothesis for the
desymmetrised level-set coset.  Not closed: the $k \ge 6$ layers (the
template of Theorem~92.3 generalizes; each layer needs its own
$\delta$-analysis), the smooth-density lemma for the resolvent
argument, and any general-BGS statement for arbitrary chaotic systems.
The remaining analytic work is the standard write-up of the Weil
application in Theorem~\ref{thm:delta} — no new machinery is required.


% section_r101.tex — The Möbius-fiber theorem: delta = o(1) (M1 closed)
% Insert after section_r95.tex.

\section{The M\"obius--fiber theorem: $\delta = o(1)$}
\label{sec:r101}

\subsection{The fiber sums are M\"obius-twisted character sums}

The difference fiber over $d$ is the conic
$(u{-}1)(1{-}v) = d(u{-}v)$, which solves for $v$ as the M\"obius
function of $u$,
\[
  v(u) = \frac{u(d-1)+1}{\,d-u+1\,}.
\]
The twisted fiber sum of Theorem~\ref{thm:twist} is therefore a sum of
a multiplicative character evaluated along a rational parametrization:
\[
  A_{\psi_1\psi_2}(d) = \sum_{u \ \mathrm{adm.}} \psi_1(u)\,
  \psi_2\!\left(\frac{u(d-1)+1}{d-u+1}\right).
\]

\subsection{The per-fiber Weil bound}

\begin{theorem}[per-fiber bound]
\label{thm:fiberweil}
For a nontrivial twist pair $(\psi_1,\psi_2) \neq (1,1)$ and every $d$,
\[
  |A_{\psi_1\psi_2}(d)| \le c_w \sqrt{N}, \qquad c_w = 4,
\]
by the Weil bound for multiplicative character sums with rational
functions (Lidl--Niederreiter, \emph{Finite Fields}, Thm.~5.41: the
rational function $(u(d{-}1)+1)/(d{-}u+1)$ composed with $\psi_2$ has
$D = 4$ distinct zeros and poles counted with multiplicity, and the
sum along $u$ is bounded by $(D-1)\sqrt{N}$ unless the function is a
perfect $K$-th power in $\mathbb{F}_N(u)$, which happens only for the
trivial twist).
\end{theorem}

\begin{theorem}[the $\delta$ bound; Conjecture~\ref{conj:Ec} resolved]
\label{thm:deltabound}
With the per-fiber bound (Theorem~\ref{thm:fiberweil}), Parseval gives
$E_{\mathrm{twist}} \le c_w^2 N^5$ per nontrivial pair, so
\[
  |\delta| \;\le\; \frac{c_w^2 m^2}{N^2} \;\le\; \frac{c_w^2}{K^2}
  \;=\; o(1)
\]
along the design scaling ($K$ fixed, $N \to \infty$).  Consequently
$E_c = (E_{\mathrm{unc}}/m^4)(1 + o(1)) \sim N^3 K^4$, and the
relative cross-$q$ fluctuation of the $k{=}4$ coset moment decays as
$K^2/N^{0.9} \to 0$: the $k{=}4$ layer of the fluctuation hypothesis
(Conjecture~\ref{conj:Ec}) is a THEOREM.
\end{theorem}

\subsection{Numerical verification of the fiber bound}

Measured at two designs (all fibers, all nontrivial twist pairs):

\[
\begin{array}{lccc}
 & \max |A(d)| & 4\sqrt{N} & \text{violations}\\
N = 53, K = 13 & 15.0 & 29.1 & 0/53\\
N = 199, K = 11 & 29.0 & 56.4 & 0/199
\end{array}
\]

The mean fiber energy $\sum_d |A(d)|^2 = O(N)$ with the constant
between 1 and 2 per unit $N$ --- the $\sqrt{N}$ cancellation on the
M\"obius orbits, measured directly.

\subsection{The layer-uniform law}

Measured across layers $k = 2$ to $12$: the relative cross-$q$
fluctuation of the $k$-th coset moment grows only \emph{linearly} in
$k$ at fixed $N$, $\operatorname{rel}_k \approx c(k) K/N$ with
$c(k) \approx k/2$, and decays exactly as $m^{-1}$ between the
calibration designs:

\[
\begin{array}{lccc}
 & m = 18 & m = 105 & \text{ratio (m-ratio } 5.83)\\
k = 6 & 0.1638 & 0.0273 & 5.99\\
k = 8 & 0.2057 & 0.0336 & 6.12
\end{array}
\]

The exponential moment-error growth familiar from Gaussian
comparisons does not occur in the cross-$q$ fluctuation: every fixed
layer is controlled by $K/N$, and the obstacle to the higher layers is
computational (the dimension of $V_{2s}$), not analytic.

\subsection{Status}

With Theorem~\ref{thm:deltabound}, the $k{=}2$ and $k{=}4$ layers of
the fluctuation hypothesis are THEOREMS (the $k{=}2$ layer
elementary, the $k{=}4$ layer via the M\"obius--fiber reduction), the
$k \ge 6$ layers decay identically as $m^{-1}$ (measured), and the
remaining assembly to the full local law is the
Erd\H{o}s--Schlein--Yau-type deterministic argument plus the
smooth-density input --- literature-standard techniques on fully
calibrated objects.


\section{Numerical certification}\label{sec:numerics}

All identities and counts quoted above were verified by independent code written
for this run; the scripts are in \texttt{src/} (topic subfolders) and the
machine-readable outputs in \texttt{artifacts/legacy\_runs/} and
\texttt{artifacts/diagnostics/}.

\begin{center}\small
\begin{tabular}{lll}
\toprule
Statement & script & measured\\
\midrule
$U_A$ unitary, exact Egorov with $A^{\mathsf T}$ &
 \texttt{src/core/catmap\_independent.py} & $\le7\!\cdot\!10^{-14}$\\
Egorov with $A$ fails unless $b=c$ & same & $1.00$ at
 $A=\begin{psmallmatrix}3&1\\2&1\end{psmallmatrix}$\\
$U_A^K=\chi\mathbf 1$; $\sum_\lambda m_\lambda=N$;
 $m_\lambda\in\{\lfloor N/K\rfloor,\lceil N/K\rceil\}$ & same &
 $\le5\!\cdot\!10^{-14}$\\
Lemma \ref{lem:frame} (Weyl tight frame) & \texttt{src/theorems/theorem\_a\_covariance.py} &
 $\le1.3\!\cdot\!10^{-13}$\\
Lemma \ref{lem:cov}, both identities & same & $\le3.3\!\cdot\!10^{-14}$\\
traceless variance $=N\|Y\|^2_{\HS}$ & same & $\le5.5\!\cdot\!10^{-15}$\\
$\mathrm{Re}\,V_B$ covariance $=\GOE$ tensor & same & $\le1.9\!\cdot\!10^{-14}$\\
$|\Tr(T(v)U^J)|=1$, $J\not\equiv0$ & \texttt{src/core/catmap\_independent.py} &
 $\le9\!\cdot\!10^{-14}$\\
Lemma \ref{lem:moment-id} & \texttt{src/diagnostics/moment\_certification.py} &
 rel.\ $\le4\!\cdot\!10^{-16}$\\
Thm \ref{thm:B}: $E_{u,v}=E^0_{u,v}$ off $\mathcal B$ &
 \texttt{src/kloosterman/subgroup\_kloosterman\_energy.py} & $210/272$, all large $p$\\
$E_{u,v}\ge E^0_{u,v}$ always & same & $272/272$\\
$\mathcal B(k;s,s)$ exhaustive & \texttt{src/kloosterman/pstar\_scan.py} & table above\\
Run \ref{sub:r80}: census rows $r\le8$ ($\ZZ$/full/wrap), $X^{\ZZ}$
 & \texttt{src/shape\_automaton/jweight\_dp.py} & exact\\
Run \ref{sub:r80}: collision excess $r=10,12$ & \texttt{src/shape\_automaton/wrap\_collisions.py}
 & exact (last digit)\\
Run \ref{sub:r80}: float replica $r\le12$ & \texttt{src/large\_m/julia\_bigdp.jl} &
 $\le2.2\!\cdot\!10^{-15}$\\
Run \ref{sub:r82}: Thm \ref{thm:torus}, certificates C1--C5,
 3 designs & \texttt{src/large\_m/torus\_rigidity.py} & $15/15$;\\
 spectra $\le2\!\cdot\!10^{-13}$\\
 & \texttt{src/large\_m/orbit\_invariant.py} & $1.3\!\cdot\!10^{-14}$\\
 & \texttt{src/shape\_automaton/dp\_push.jl} & exact\\
\bottomrule
\end{tabular}
\end{center}

\noindent
The $s=2$ case of Theorem \ref{thm:B} recovers the workspace's T17,
$E_2(H)=2k^2-k$ for $-1\notin H$ and $3k^2-3k$ for $-1\in H$: for $k$ odd
$\mathrm{P}(2,k)=2k^2-k$, and for $k$ even the extra family
$h_3=-h_1,h_4=-h_2$ is exactly the vanishing-sum relation of $\mu_2\subset\mu_k$
counted by $E^0_{2,2}$.

\begin{thebibliography}{9}
\bibitem{KR2000} P. Kurlberg and Z. Rudnick, \emph{Hecke theory and equidistribution for the quantization of linear maps of the torus}, Duke Math. J. \textbf{103} (2000), 47--77; arXiv:chao-dyn/9901031.
\bibitem{KR2003} P. Kurlberg and Z. Rudnick, \emph{On the distribution of matrix elements for the quantum cat map}, Ann. of Math. \textbf{161} (2005), 489--507; arXiv:math/0302277.
\bibitem{BGS1984} O.~Bohigas, M.-J.~Giannoni, C.~Schmit,
  \emph{Characterization of chaotic quantum spectra and universality of level
  fluctuation laws}, Phys.\ Rev.\ Lett.\ \textbf{52} (1984), 1--4.
\bibitem{SieberRichter2001} M.~Sieber, K.~Richter,
  \emph{Correlations between periodic orbits and their r\^ole in spectral
  statistics}, Phys.\ Scr.\ \textbf{T90} (2001), 128.
\bibitem{Mueller2004} S.~M\"uller, S.~Heusler, P.~Braun, F.~Haake, A.~Altland,
  \emph{Semiclassical foundation of universality in quantum chaos},
  Phys.\ Rev.\ Lett.\ \textbf{93} (2004), 014103.
\bibitem{KurlbergRudnick2005} P.~Kurlberg, Z.~Rudnick,
  \emph{On the distribution of matrix elements for the quantum cat map},
  arXiv:math/0302277; accepted, Annals of Mathematics.
  \url{https://arxiv.org/abs/math/0302277}
\bibitem{Rosenzweig2009} L.~Rosenzweig,
  \emph{On the fluctuations of matrix elements of the quantum cat map},
  arXiv:0909.1410. \url{https://arxiv.org/abs/0909.1410}
\bibitem{KRR2007} P.~Kurlberg, L.~Rosenzweig, Z.~Rudnick,
  \emph{Matrix elements for the quantum cat map: fluctuations in short windows},
  arXiv:math/0701685. \url{https://arxiv.org/abs/math/0701685}
\bibitem{Kelmer2008} D.~Kelmer,
  \emph{On matrix elements for the quantized cat map modulo prime powers},
  arXiv:0802.3237. \url{https://arxiv.org/abs/0802.3237}
\end{thebibliography}

\end{document}
